% Part IV: Formal Proofs
% This file is \include'd from main.tex

%%%%%%%%%%%%%%%%%%%%%%%%%%%%%%%%%%%%%%%%%%%%%%%%%%%%%%%%%%%%%%%%%%%%%%%%
%%  Chapter 18 — Proof of Equivalence
%%%%%%%%%%%%%%%%%%%%%%%%%%%%%%%%%%%%%%%%%%%%%%%%%%%%%%%%%%%%%%%%%%%%%%%%
\chapter{Equivalence of Structured and Component-Wise Views}
\label{ch:proof-equiv}

This chapter contains the full proof that the category of structured
(4-tuple) judgments $\mathbf{Judg}_4$ is equivalent to the category of
component-wise judgments $\mathbf{Judg}_8$.  We construct the relevant
categories, the flattening and structuring functors between them, and
verify all required natural isomorphisms.

%% ── §18.1  The Category Judg_8 ──────────────────────────────────────
\section{The category $\mathbf{Judg}_8$}
\label{sec:judg8}

\begin{definition}[Objects of $\mathbf{Judg}_8$]
\label{def:judg8-obj}
An object of $\mathbf{Judg}_8$ is a component tuple
\[
  J = (c,\,\varphi,\,A,\,E,\,O,\,B,\,T,\,\Pi)
\]
where:
\begin{enumerate}[label=(\roman*),nosep]
\item $c \in \ob(\mathcal{C})$ is a coordinate in a small category $\mathcal{C}$
  equipped with a Grothendieck topology~$\mathcal{J}$;
\item $\varphi \in \PSh(\mathcal{C})(c)$ is a local section of a
  presheaf $\mathcal{F}$ at~$c$;
\item $A$ is a type in a dependent type theory $\mathcal{T}$;
\item $E$ is a term such that $\Gamma \vdash E : A$ for some context~$\Gamma$;
\item $O = \{U_i \to X\}_{i \in I}$ is a covering sieve (obligation cover);
\item $B \in \check{H}^1(\mathcal{U}, \mathcal{F})$ is a \v{C}ech
  cohomology class (obstruction);
\item $T \in \Grade$ is a grade;
\item $\Pi$ is a provenance trail (a finite record of evidence sources
  and derivation steps).
\end{enumerate}
These satisfy the \emph{coherence conditions}:
\begin{enumerate}[label=(C\arabic*),nosep]
\item The cover~$O$ refines the topology~$\mathcal{J}$:
  $O \in \mathcal{J}(X)$.
\item The class~$B$ is computed from the presheaf~$\mathcal{F}$
  restricted to the cover~$O$.
\item The grade~$T$ is the evaluation of the trail:
  $T = \mu(\Pi)$, where $\mu$ is the grade-evaluation map.
\item The typing judgment $\Gamma \vdash E : A$ is derivable in~$\mathcal{T}$.
\end{enumerate}
\end{definition}

\begin{definition}[Morphisms of $\mathbf{Judg}_8$]
\label{def:judg8-mor}
A morphism $f : J \to J'$ in $\mathbf{Judg}_8$ is a component tuple
\[
  f = (f_c,\, f_\varphi,\, f_A,\, f_E,\, f_O,\, f_B,\, f_T,\, f_\Pi)
\]
where:
\begin{enumerate}[label=(\roman*),nosep]
\item $f_c : c \to c'$ is a morphism in $\mathcal{C}$;
\item $f_\varphi : \mathcal{F}(c') \to \mathcal{F}(c)$ is the
  restriction along $f_c$, satisfying $f_\varphi(\varphi') = \varphi$;
\item $f_A : A \to A'$ is a type morphism (a term $f_A : A \to A'$
  in context~$\Gamma$);
\item $f_E$ witnesses the transport:
  $\Gamma \vdash f_E : f_A(E) =_{A'} E'$;
\item $f_O$ is a refinement of covers: every $U_i' \in O'$ is covered by
  members of~$O$;
\item $f_B : \check{H}^1(\mathcal{U},\mathcal{F}) \to
  \check{H}^1(\mathcal{U}',\mathcal{F}')$ is the induced map on
  cohomology, satisfying $f_B(B) = B'$;
\item $f_T$ witnesses $T' \leq T$ (grade is non-increasing under morphisms);
\item $f_\Pi$ is an extension of trails: $\Pi' = f_\Pi(\Pi)$.
\end{enumerate}
Composition is componentwise. The identity morphism is the tuple of
identity maps.
\end{definition}

\begin{lemma}
\label{lem:judg8-cat}
$\mathbf{Judg}_8$ is a category.
\end{lemma}

\begin{proof}
We verify the category axioms.
\emph{Identity:} For each object~$J$, the tuple
$(\id_c, \id_{\mathcal{F}(c)}, \id_A, \mathrm{refl}_E, \id_O, \id_B, \mathrm{refl}_T, \id_\Pi)$
satisfies all eight morphism conditions, since identity maps preserve
all structure.

\emph{Composition:} Given $f : J \to J'$ and $g : J' \to J''$, define
$g \circ f$ componentwise:
$(g_c \circ f_c,\; f_\varphi \circ g_\varphi,\; g_A \circ f_A,\;
\ldots)$.
The reversal of order for~$f_\varphi$ follows from contravariance of
presheaves.  Each coherence condition is preserved because each
component morphism preserves its respective condition.

\emph{Associativity:} Follows from associativity of composition in each
ambient category ($\mathcal{C}$, $\Set$, the type theory, etc.).

\emph{Unit laws:} The identity tuple is a two-sided unit for
componentwise composition, since identity is a unit in each factor.
\end{proof}

%% ── §18.2  The Category Judg_4 ──────────────────────────────────────
\section{The category $\mathbf{Judg}_4$}
\label{sec:judg4}

\begin{definition}[Objects of $\mathbf{Judg}_4$]
\label{def:judg4-obj}
An object of $\mathbf{Judg}_4$ is a 4-tuple
\[
  \JJ = (\STop,\; \PStr,\; \CSpec,\; \GEval)
\]
where:
\begin{enumerate}[label=(\roman*),nosep]
\item $\STop = (\mathcal{C}, \mathcal{J}, \mathcal{F}, \Omega,
  a, \gamma, \tau)$ is a \emph{Semantic Topos}:
  a Grothendieck topos $\Sh(\mathcal{C}, \mathcal{J})$ with a
  distinguished presheaf~$\mathcal{F}$, subobject classifier~$\Omega$,
  sheafification~$a$, global-sections geometric morphism~$\gamma$,
  and Lawvere--Tierney topology~$\tau$.
\item $\PStr = (\Gamma, \mathcal{T}, E, A, \mathcal{D})$ is a
  \emph{Proof Structure}: a context~$\Gamma$, a dependent type
  theory~$\mathcal{T}$ with universes, a derivation tree~$\mathcal{D}$
  witnessing $\Gamma \vdash E : A$, together with reduction rules
  ($\beta$, $\eta$) and the full apparatus of dependent types
  ($\Pi$, $\Sigma$, $\mathrm{Id}$, $W$-types).
\item $\CSpec = (\mathcal{U}, C^\bullet, H^\bullet, \partial, \mathrm{SS})$
  is a \emph{Cohomological Spectrum}: a cover~$\mathcal{U}$, the full
  \v{C}ech complex $C^\bullet(\mathcal{U}, \mathcal{F})$, the
  cohomology groups $H^n$ for all $n \geq 0$, connecting maps, and
  the \v{C}ech-to-derived spectral sequence~$\mathrm{SS}$.
\item $\GEval = (G, \opG, \otG, \gzero, \gone, \mu, \multimap, \leq)$
  is a \emph{Graded Evaluation}: a commutative unital quantale $G$
  with evaluation functor~$\mu$, residuation~$\multimap$, and the
  provenance-tracking structure.
\end{enumerate}
These satisfy the \emph{compatibility conditions}:
\begin{enumerate}[label=(K\arabic*),nosep]
\item The cover~$\mathcal{U}$ in $\CSpec$ refines the
  topology~$\mathcal{J}$ in $\STop$.
\item The presheaf~$\mathcal{F}$ in $\STop$ is the coefficient sheaf
  of the \v{C}ech complex in $\CSpec$.
\item The type~$A$ in $\PStr$ classifies sections of~$\mathcal{F}$:
  global sections of~$\mathcal{F}$ correspond to closed terms of type~$A$.
\item The evaluation $\mu : \mathcal{D} \to G$ in $\GEval$ takes
  derivation trees from~$\PStr$ to grades in~$G$.
\end{enumerate}
\end{definition}

\begin{definition}[Morphisms of $\mathbf{Judg}_4$]
\label{def:judg4-mor}
A morphism $F : \JJ \to \JJ'$ in $\mathbf{Judg}_4$ is a 4-tuple
\[
  F = (F_{\STop},\; F_{\PStr},\; F_{\CSpec},\; F_{\GEval})
\]
where:
\begin{enumerate}[label=(\roman*),nosep]
\item $F_{\STop} : \STop \to \STop'$ is a geometric morphism of
  topoi, i.e.\ an adjoint pair $(f^* \dashv f_*)$ with $f^*$
  left-exact, mapping $\mathcal{F} \mapsto \mathcal{F}'$;
\item $F_{\PStr} : \PStr \to \PStr'$ is a context morphism: a
  substitution $\sigma : \Gamma \to \Gamma'$ such that
  $E'[\sigma] \equiv E$ and $A'[\sigma] \equiv A$ (up to
  definitional equality);
\item $F_{\CSpec} : \CSpec \to \CSpec'$ is a chain map
  $C^\bullet(\mathcal{U}, \mathcal{F}) \to
  C^\bullet(\mathcal{U}', \mathcal{F}')$ compatible with coboundary maps;
\item $F_{\GEval} : \GEval \to \GEval'$ is a quantale homomorphism
  $\alpha : G \to G'$ preserving $\opG$, $\otG$, $\gzero$, $\gone$,
  and satisfying $\alpha \circ \mu = \mu' \circ F_{\PStr}$.
\end{enumerate}
These satisfy: $F_{\CSpec}$ is the chain map induced by $F_{\STop}$
on \v{C}ech complexes, and $F_{\GEval}$ is compatible with
$F_{\PStr}$ via the evaluation functors.

Composition is componentwise (composing geometric morphisms, composing
substitutions, composing chain maps, composing quantale homomorphisms).
The identity morphism is the tuple of identity maps.
\end{definition}

\begin{lemma}
\label{lem:judg4-cat}
$\mathbf{Judg}_4$ is a category.
\end{lemma}

\begin{proof}
Composition of geometric morphisms is a geometric morphism
\cite{johnstone2002}.  Composition of substitutions is a substitution.
Composition of chain maps is a chain map.  Composition of quantale
homomorphisms is a quantale homomorphism.  Compatibility conditions
(K1)--(K4) are preserved because each component functor preserves the
relevant structure.  Associativity and unit laws follow from those in
each ambient category.
\end{proof}

%% ── §18.3  The Forgetful Functor ────────────────────────────────────
\section{The forgetful functor $U : \mathbf{Judg}_4 \to \mathbf{Judg}_8$}
\label{sec:forgetful}

\begin{construction}[The functor $U$]
\label{con:forgetful}
Define $U$ as follows.

\emph{On objects.} Given $\JJ = (\STop, \PStr, \CSpec, \GEval)$, set
\[
  U(\JJ) = (c,\, \varphi,\, A,\, E,\, O,\, B,\, T,\, \Pi)
\]
where:
\begin{itemize}[nosep]
\item $c$ is any chosen base object of $\mathcal{C}$ (we fix a global
  choice of base point for the site; for pointed sites this is
  canonical);
\item $\varphi = \mathcal{F}(c)$ is the stalk of the presheaf at~$c$;
\item $(A, E)$ are the type and evidence term from $\PStr$;
\item $O = \mathcal{U}$ is the cover from $\CSpec$;
\item $B = [z] \in H^1(\mathcal{U}, \mathcal{F})$ is the first
  cohomology class from $\CSpec$;
\item $T = \mu(\mathcal{D})$ is the grade of the derivation;
\item $\Pi$ is the provenance record extracted from $\GEval$.
\end{itemize}

\emph{On morphisms.} Given $F = (F_{\STop}, F_{\PStr}, F_{\CSpec},
F_{\GEval})$, extract the underlying component maps:
\begin{itemize}[nosep]
\item $f_c$ is the image of the base point under the inverse image
  functor $f^*$;
\item $f_\varphi$ is the map on stalks induced by $F_{\STop}$;
\item $f_A$ is the type map induced by the substitution in $F_{\PStr}$;
\item $f_E$ is the transport witness from $F_{\PStr}$;
\item $f_O$ is the refinement of covers induced by $F_{\CSpec}$;
\item $f_B$ is the map $H^1(F_{\CSpec})$ on first cohomology;
\item $f_T$ is the grade map from $F_{\GEval}$;
\item $f_\Pi$ is the provenance extension.
\end{itemize}
\end{construction}

\begin{proposition}
\label{prop:U-functor}
$U : \mathbf{Judg}_4 \to \mathbf{Judg}_8$ is a functor.
\end{proposition}

\begin{proof}
\emph{Preservation of identities.}
$U(\id_\JJ) = (\id_c, \id_\varphi, \id_A, \mathrm{refl}_E,
\id_O, \id_B, \mathrm{refl}_T, \id_\Pi)$,
which is the identity in $\mathbf{Judg}_8$ by
\cref{def:judg8-mor}.

\emph{Preservation of composition.}
Let $F : \JJ \to \JJ'$ and $G : \JJ' \to \JJ''$.  Then
$U(G \circ F) = (g_c \circ f_c,\; f_\varphi \circ g_\varphi,\;
\ldots) = U(G) \circ U(F)$,
where the reversal in the second component follows from
contravariance of presheaf restriction and the fact that the inverse
image functor of the composite geometric morphism is the composite of
the inverse image functors.  Each remaining component is computed by
the functoriality of the relevant extraction (chain maps compose,
quantale homomorphisms compose, etc.).
\end{proof}

%% ── §18.4  The Structuring Functor ───────────────────────────────────
\section{The structuring functor $R : \mathbf{Judg}_8 \to \mathbf{Judg}_4$}
\label{sec:enrichment}

\begin{construction}[The functor $R$]
\label{con:enrichment}
Given a component tuple $J = (c, \varphi, A, E, O, B, T, \Pi)$ satisfying
the coherence conditions (C1)--(C4), we construct a 4-tuple
$R(J) = (\STop_J, \PStr_J, \CSpec_J, \GEval_J)$ as follows.

\paragraph{Semantic Topos $\STop_J$.}
The site $(\mathcal{C}, \mathcal{J})$ is given as part of the data.
The presheaf $\mathcal{F}$ is defined by $\varphi \in \mathcal{F}(c)$.
We set $\STop_J$ to be the \emph{smallest sub-topos} of
$\PSh(\mathcal{C})$ that:
\begin{enumerate}[label=(\alph*),nosep]
\item contains the presheaf $\mathcal{F}$,
\item has a topology that refines $\mathcal{J}$,
\item is closed under the internal logic operations
  (finite limits, exponentials, subobject classifier).
\end{enumerate}
Concretely, $\STop_J = \Sh(\mathcal{C}, \mathcal{J}_{\mathrm{min}})$
where $\mathcal{J}_{\mathrm{min}}$ is the coarsest topology on
$\mathcal{C}$ refining~$\mathcal{J}$ for which $\mathcal{F}$ is a
sheaf.  Such a topology exists by the intersection property of
Grothendieck topologies \cite{johnstone2002}.
The remaining components ($\Omega, a, \gamma, \tau$) are canonical:
$\Omega$ is the subobject classifier of $\Sh(\mathcal{C},
\mathcal{J}_{\mathrm{min}})$, $a$ is sheafification,
$\gamma$ is the global-sections geometric morphism to $\Set$,
and $\tau$ is the Lawvere--Tierney topology corresponding to
$\mathcal{J}_{\mathrm{min}}$.

\paragraph{Proof Structure $\PStr_J$.}
From $(A, E)$ and the derivability condition (C4), there exists a
context $\Gamma$ and a derivation tree $\mathcal{D}$ witnessing
$\Gamma \vdash E : A$.  We set $\PStr_J$ to be the
\emph{minimal type theory} containing:
\begin{enumerate}[label=(\alph*),nosep]
\item the types appearing in $\Gamma$ and $A$,
\item the term constructors used in~$E$,
\item identity types $\mathrm{Id}$ for all types in the theory,
\item closure under $\Pi$-types, $\Sigma$-types, and the universe
  $\Type_0$ containing~$A$.
\end{enumerate}
This is well-defined: the intersection of all type theories containing
the given data is again a type theory (closed under the type formation
rules).

\paragraph{Cohomological Spectrum $\CSpec_J$.}
The cover $O = \mathcal{U}$ and the presheaf $\mathcal{F}$ determine
the full \v{C}ech complex
\[
  C^0(\mathcal{U}, \mathcal{F}) \xrightarrow{\delta^0}
  C^1(\mathcal{U}, \mathcal{F}) \xrightarrow{\delta^1}
  C^2(\mathcal{U}, \mathcal{F}) \xrightarrow{\delta^2} \cdots
\]
We take $\CSpec_J$ to be this complex together with all its
cohomology groups $H^n = \ker \delta^n / \mathrm{im}\, \delta^{n-1}$.
The obstruction class~$B$ is recovered as $B = [z] \in H^1$.
The \v{C}ech-to-derived spectral sequence is the standard one
\cite{weibel1994}.

\paragraph{Graded Evaluation $\GEval_J$.}
The grade $T$ lives in a grade set $G$.  We equip $G$ with the
\emph{canonical quantale structure}: define $\opG = \max$,
$\otG = \min$ (the ``bottleneck'' semiring on a totally ordered set),
$\gzero = \bot$, $\gone = \top$.  The evaluation map $\mu$ is defined
by $\mu(\Pi) = T$ (extending to all derivation trees by the
recursive grade computation).
The residuation is $a \multimap b = \begin{cases} \gone & \text{if }
a \leq b, \\ b & \text{otherwise.}\end{cases}$
\end{construction}

\begin{proposition}
\label{prop:R-objects}
For every object $J$ of $\mathbf{Judg}_8$, the tuple $R(J)$ is a
well-defined object of $\mathbf{Judg}_4$.
\end{proposition}

\begin{proof}
We verify the compatibility conditions (K1)--(K4) of
\cref{def:judg4-obj}.

(K1): The cover $\mathcal{U}$ in $\CSpec_J$ is exactly the obligation
cover~$O$, which refines $\mathcal{J}$ by coherence condition (C1),
and $\mathcal{J}_{\mathrm{min}}$ refines~$\mathcal{J}$, so $\mathcal{U}$
refines $\mathcal{J}_{\mathrm{min}}$.

(K2): The presheaf $\mathcal{F}$ used in $\CSpec_J$ is the same
presheaf present in $\STop_J$ by construction.

(K3): The type~$A$ in $\PStr_J$ classifies sections of $\mathcal{F}$
because we have included in the minimal type theory exactly the types
needed to represent sections, and the evidence term~$E$ of type~$A$ is
a witness of a section at coordinate~$c$.

(K4): The evaluation $\mu(\mathcal{D}) = T$ by coherence condition (C3)
and the definition of $\mu$ in $\GEval_J$.
\end{proof}

\begin{construction}[$R$ on morphisms]
\label{con:R-morphisms}
Given a morphism $f = (f_c, f_\varphi, f_A, f_E, f_O, f_B, f_T,
f_\Pi) : J \to J'$ in $\mathbf{Judg}_8$, we define
$R(f) = (F_{\STop}, F_{\PStr}, F_{\CSpec}, F_{\GEval})$ as follows:
\begin{enumerate}[label=(\roman*),nosep]
\item $F_{\STop}$ is the geometric morphism induced by the site
  morphism $f_c : \mathcal{C} \to \mathcal{C}'$ (extended to a
  continuous functor between sites); the inverse image
  $f^*$ sends $\mathcal{F}'$ to $\mathcal{F}$ via $f_\varphi$.
\item $F_{\PStr}$ is the substitution $\sigma$ induced by
  $(f_A, f_E)$: define $\sigma(x) = f_A(x)$ for type variables and
  $f_E$ provides the transport witness.
\item $F_{\CSpec}$ is the chain map induced by $f_O$ (refinement of
  covers) and $f_\varphi$ (the presheaf map): on $n$-cochains,
  the map is the composite of restriction along the refinement
  and the presheaf map.
\item $F_{\GEval}$ is the quantale homomorphism
  $\alpha(g) = f_T(g)$ for $g \in G$, which preserves
  $\opG, \otG$ because $f_T$ is monotone and the quantale
  operations are order-theoretic.
\end{enumerate}
\end{construction}

\begin{proposition}
\label{prop:R-functor}
$R : \mathbf{Judg}_8 \to \mathbf{Judg}_4$ is a functor.
\end{proposition}

\begin{proof}
\emph{Preservation of identities.}  $R(\id_J) = (\id_{\STop_J},
\id_{\PStr_J}, \id_{\CSpec_J}, \id_{\GEval_J})$, since the identity
site morphism induces the identity geometric morphism, the identity
substitution is the identity on contexts, the identity refinement
induces the identity chain map, and the identity on grades is a
quantale homomorphism.

\emph{Preservation of composition.}  Let $f : J \to J'$ and
$g : J' \to J''$.  Each component of $R(g \circ f)$ is computed from
the composite $(g_c \circ f_c, \ldots)$.  The induced geometric
morphism of a composite site morphism is the composite of the
induced geometric morphisms \cite{johnstone2002}, so
$(R(g) \circ R(f))_{\STop} = R(g \circ f)_{\STop}$.  Similarly for
substitutions (composition of substitutions), chain maps (composition
of chain maps), and quantale homomorphisms.
\end{proof}

%% ── §18.5  U ∘ R ≅ id ──────────────────────────────────────────────
\section{The isomorphism $U \circ R \cong \id_{\mathbf{Judg}_8}$}
\label{sec:UR}

\begin{theorem}
\label{thm:UR}
There is a natural isomorphism $\eta : \id_{\mathbf{Judg}_8}
\Rightarrow U \circ R$.
\end{theorem}

\begin{proof}
Let $J = (c, \varphi, A, E, O, B, T, \Pi)$ be an object of
$\mathbf{Judg}_8$.  We must exhibit an isomorphism
$\eta_J : J \xrightarrow{\sim} (U \circ R)(J)$ and verify
naturality.

\emph{The component $\eta_J$.}  By construction,
$(U \circ R)(J) = U(\STop_J, \PStr_J, \CSpec_J, \GEval_J)$.
The forgetful functor~$U$ extracts:
\begin{itemize}[nosep]
\item the base point of $\STop_J$: this is $c$, since $\STop_J$
  was built from the site containing~$c$;
\item the stalk $\mathcal{F}(c)$: this contains $\varphi$ by
  construction of $\STop_J$;
\item $(A, E)$: unchanged, since $\PStr_J$ was the minimal theory
  containing $(A, E)$;
\item $O$: unchanged, since $\CSpec_J$ was built from the cover~$O$;
\item $B$: unchanged, since $H^1$ of the \v{C}ech complex built
  from $(O, \mathcal{F})$ recovers~$B$;
\item $T = \mu(\Pi)$: unchanged by (C3) and the definition of
  $\GEval_J$;
\item $\Pi$: unchanged.
\end{itemize}
Thus $(U \circ R)(J) = J$, and $\eta_J = \id_J$.

\emph{Naturality.}  For any morphism $f : J \to J'$ in
$\mathbf{Judg}_8$, the naturality square
\[
\begin{tikzcd}
  J \ar[r, "\eta_J"] \ar[d, "f"'] &
  (U \circ R)(J) \ar[d, "(U \circ R)(f)"] \\
  J' \ar[r, "\eta_{J'}"'] & (U \circ R)(J')
\end{tikzcd}
\]
commutes trivially, since $\eta_J = \id_J$ and
$\eta_{J'} = \id_{J'}$, so both paths equal~$f$ (using
$(U \circ R)(f) = f$, which follows from the fact that $U$ and $R$
act as identity on the underlying component data).
\end{proof}

%% ── §18.6  R ∘ U ≅ id ──────────────────────────────────────────────
\section{The isomorphism $R \circ U \cong \id_{\mathbf{Judg}_4}$}
\label{sec:RU}

This direction is more substantive: we must show that the canonical
enrichment of the extracted component data is isomorphic to the given
4-tuple.  The key insight is that the ``extra'' structure in a 4-tuple
object is \emph{determined up to canonical isomorphism} by its
underlying components.

\begin{theorem}
\label{thm:RU}
There is a natural isomorphism $\varepsilon : R \circ U
\Rightarrow \id_{\mathbf{Judg}_4}$.
\end{theorem}

\begin{proof}
Let $\JJ = (\STop, \PStr, \CSpec, \GEval)$ be an object of
$\mathbf{Judg}_4$.  We must exhibit an isomorphism
$\varepsilon_\JJ : (R \circ U)(\JJ) \xrightarrow{\sim} \JJ$.

\emph{Step 1: The topos component.}
$(R \circ U)(\JJ)$ has topos $\STop_{U(\JJ)}$, which is the smallest
sub-topos of $\PSh(\mathcal{C})$ containing the presheaf~$\mathcal{F}$
extracted from~$\STop$.  Now $\STop$ itself is a sub-topos of
$\PSh(\mathcal{C})$ containing~$\mathcal{F}$, and $\STop_{U(\JJ)}$ is
the \emph{smallest} such.  Therefore there is a canonical geometric
inclusion
\[
  \iota_{\STop} : \STop_{U(\JJ)} \hookrightarrow \STop.
\]
For this to be an isomorphism, we need $\STop$ to be generated by
$\mathcal{F}$ and the site data---that is, $\STop$ carries no
``extra'' sheaves beyond those determined by the site and the
presheaf.

We claim this holds for the objects of $\mathbf{Judg}_4$: by
compatibility condition (K2), the presheaf $\mathcal{F}$ in $\STop$
is the same as the coefficient sheaf in $\CSpec$, and by (K1), the
topology is determined by the covering data.  Since $\STop$ was
defined as the topos generated by the site and $\mathcal{F}$ (with
the Lawvere--Tierney topology $\tau$ being the unique one making
$\mathcal{F}$ a sheaf for $\mathcal{J}$), we have
$\STop_{U(\JJ)} \cong \STop$.  The isomorphism is the canonical
geometric morphism $\varepsilon_{\STop}$.

\emph{Step 2: The type-theory component.}
$(R \circ U)(\JJ)$ has proof structure $\PStr_{U(\JJ)}$, which is the
minimal type theory containing $(A, E, \Gamma)$ extracted
from~$\PStr$.  Now $\PStr$ contains $(A, E, \Gamma)$ and possibly
additional types and terms.  However, the minimality of $\PStr$ as a
proof structure for the judgment $\JJ$ (every component of $\PStr$ is
used in the derivation~$\mathcal{D}$ or required by the closure
conditions) ensures that $\PStr_{U(\JJ)}$ and $\PStr$ have the same
types and terms up to definitional equality.  Thus
$\varepsilon_{\PStr} : \PStr_{U(\JJ)} \xrightarrow{\sim} \PStr$
is an isomorphism of type theories.

\emph{Step 3: The cohomological component.}
$(R \circ U)(\JJ)$ has spectrum $\CSpec_{U(\JJ)}$, which is the
\v{C}ech complex of $(\mathcal{U}, \mathcal{F})$ with all cohomology
groups.  Since $\CSpec$ is \emph{defined} as this \v{C}ech complex
(compatibility conditions (K1)--(K2) ensure the cover and sheaf
agree), we have $\CSpec_{U(\JJ)} = \CSpec$ on the nose.
Set $\varepsilon_{\CSpec} = \id$.

\emph{Step 4: The grade component.}
$(R \circ U)(\JJ)$ has evaluation $\GEval_{U(\JJ)}$, which is the
canonical quantale on $G$ with $\mu(\Pi) = T$.  The given $\GEval$
has the same underlying set $G$ and the same evaluation $\mu$
(by (K4)).  The quantale structure is determined by the total
order on $G$ (since $\opG = \max$ and $\otG = \min$ are the only
idempotent quantale operations on a totally ordered set that have
$\gone = \top$ and $\gzero = \bot$).  Thus
$\varepsilon_{\GEval} : \GEval_{U(\JJ)} \xrightarrow{\sim} \GEval$
is an isomorphism.

\emph{Naturality.}  Let $F : \JJ \to \JJ'$ be a morphism in
$\mathbf{Judg}_4$.  We must show that the diagram
\[
\begin{tikzcd}
  (R \circ U)(\JJ) \ar[r, "\varepsilon_\JJ"]
    \ar[d, "(R \circ U)(F)"'] &
  \JJ \ar[d, "F"] \\
  (R \circ U)(\JJ') \ar[r, "\varepsilon_{\JJ'}"'] &
  \JJ'
\end{tikzcd}
\]
commutes.  Each component of $\varepsilon$ is either an identity or a
canonical inclusion/isomorphism.  Geometric morphisms compose
functorially, so $F_{\STop} \circ \varepsilon_{\STop} =
\varepsilon_{\STop'} \circ (R \circ U)(F)_{\STop}$ by
functoriality of the sub-topos inclusion.  The type-theory and
grade components are analogous: canonical isomorphisms commute with
structure-preserving maps.  The cohomological component is strictly
an identity, so the square commutes trivially.
\end{proof}

%% ── §18.7  Main Theorem and Corollaries ─────────────────────────────
\section{The equivalence theorem}
\label{sec:equiv-thm}

\begin{theorem}[Categorical equivalence of views]
\label{thm:equiv}
The functors $U : \mathbf{Judg}_4 \to \mathbf{Judg}_8$ and
$R : \mathbf{Judg}_8 \to \mathbf{Judg}_4$ exhibit an equivalence
of categories:
\[
  \mathbf{Judg}_4 \simeq \mathbf{Judg}_8.
\]
\end{theorem}

\begin{proof}
By \cref{thm:UR}, $U \circ R \cong \id_{\mathbf{Judg}_8}$.
By \cref{thm:RU}, $R \circ U \cong \id_{\mathbf{Judg}_4}$.
Hence $U$ and $R$ are quasi-inverse functors, establishing the
equivalence.
\end{proof}

\begin{corollary}[Transfer principle]
\label{cor:transfer}
Let $P$ be any property expressible in the internal language of
$\mathbf{Judg}_8$ (respectively $\mathbf{Judg}_4$).  Then $P$ holds
for a component tuple $J$ if and only if $P$ holds for the corresponding
4-tuple $R(J)$ (respectively, $P$ holds for a 4-tuple $\JJ$ if and
only if $P$ holds for $U(\JJ)$).
\end{corollary}

\begin{proof}
Equivalences of categories preserve and reflect all categorical
properties: limits, colimits, (co)equalizers, pullbacks, exponentials,
and internal-logic statements.  Since $U$ and $R$ form an equivalence,
any property of $J$ in $\mathbf{Judg}_8$ is transferred to $R(J)$ in
$\mathbf{Judg}_4$ via the natural isomorphism $\eta$, and conversely
via $\varepsilon$.
\end{proof}

\begin{corollary}[Faithfulness of $U$]
\label{cor:faithful}
The forgetful functor $U$ is faithful: distinct morphisms in
$\mathbf{Judg}_4$ have distinct underlying component-wise morphisms.
\end{corollary}

\begin{proof}
Any functor that is part of an equivalence is fully faithful.
Faithfulness follows immediately.
\end{proof}

\begin{remark}[Non-injectivity on objects]
\label{rem:non-inj}
Although $U$ is faithful on morphisms, it is not injective on objects
in general: distinct topoi $\STop \neq \STop'$ can project to the
same pair $(c, \varphi)$ if different Grothendieck topologies yield
the same stalks.  The equivalence resolves this through the enrichment
functor $R$, which selects the \emph{canonical} (minimal) enrichment.
Objects that project identically under $U$ are identified up to the
natural isomorphism~$\varepsilon$.
\end{remark}


%% ── Stub for remaining chapters ─────────────────────────────────────
%%%%%%%%%%%%%%%%%%%%%%%%%%%%%%%%%%%%%%%%%%%%%%%%%%%%%%%%%%%%%%%%%%%%%%%%
%%  Chapter 19 — Subsumption Theorems
%%%%%%%%%%%%%%%%%%%%%%%%%%%%%%%%%%%%%%%%%%%%%%%%%%%%%%%%%%%%%%%%%%%%%%%%
\chapter{Subsumption Theorems}
\label{ch:subsumption}

This chapter proves that the JuGeo 4-tuple framework properly
subsumes eight established formalisms.  For each, we give the precise
definition of the subsumed framework, construct a faithful embedding
into $\mathbf{Judg}_4$, prove that the embedding preserves the
relevant structure, and exhibit a concrete example.

\medskip

Throughout, we write $\mathbf{2} = \{\gzero, \gone\}$ for the
two-element Boolean grade semiring with $\opG = \lor$ and
$\otG = \land$, and we write $\mathbf{1}$ for the terminal
category (single object, identity morphism only).

%% ── §19.1  Montague Grammar ────────────────────────────────────────
\section{JuGeo subsumes Montague Grammar}
\label{sec:sub-montague}

\begin{definition}[Montague IL]
\label{def:montague}
Montague's Intensional Logic (IL) is a simply-typed $\lambda$-calculus
with types generated by base types $e$ (entities) and $t$ (truth
values), function types $\alpha \to \beta$, and an intension operator
$\hat{\phantom{x}}$.  A \emph{Montague model} is a triple
$(W, D, \llbracket\cdot\rrbracket)$ where $W$ is a set of possible
worlds, $D$ is a domain of entities, and
$\llbracket\cdot\rrbracket : \mathrm{Term} \to (W \to D_\alpha)$
is an interpretation function.
\end{definition}

\begin{construction}[Embedding $\iota_{\mathrm{MG}} :
  \mathbf{Mont} \to \mathbf{Judg}_4$]
\label{con:montague-embed}
Given a Montague model $(W, D, \llbracket\cdot\rrbracket)$ and a
typed term $M : \alpha$, define:
\begin{itemize}[nosep]
\item $\STop$: Let $\mathcal{C} = \mathcal{O}(W)$ be the poset of
  open subsets of~$W$ (with the usual topology).  The presheaf
  $\mathcal{F}(U) = \{f : U \to D_\alpha\}$ assigns to each open set
  the set of $\alpha$-typed functions on it.  Sheafification gives the
  possible-worlds topos $\Sh(W)$.

\item $\PStr$: The simply-typed $\lambda$-calculus with types built
  from $e$ and $t$ embeds into dependent type theory by setting
  $\Gamma = (w : W)$ and $A_\alpha = D_\alpha$.  The term $M$ becomes
  an evidence term $E = \lambda w.\, \llbracket M \rrbracket(w)$.

\item $\CSpec$: Since $\Sh(W)$ is a topos of sheaves on a
  topological space, and the presheaf $\mathcal{F}$ is already a
  sheaf (function sheaves are sheaves), we have $H^n = 0$ for all
  $n \geq 1$ on any open cover.  Set $\CSpec$ to be the trivial
  \v{C}ech complex.

\item $\GEval$: Set $G = \mathbf{2}$ (binary grades).  A sentence
  $\varphi$ has grade $\gone$ if $\llbracket\varphi\rrbracket(w) =
  \mathrm{true}$ and $\gzero$ otherwise.  The evaluation functor is
  $\mu(\mathcal{D}) = \gone$ if the derivation is valid, $\gzero$
  otherwise.
\end{itemize}
\end{construction}

\begin{theorem}
\label{thm:montague}
The embedding $\iota_{\mathrm{MG}}$ is faithful and preserves
truth conditions.
\end{theorem}

\begin{proof}
\emph{Faithfulness.}  Let $f, g : M \to M'$ be distinct morphisms in
$\mathbf{Mont}$ (i.e., distinct type-preserving maps between
Montague models).  The embedding maps these to geometric morphisms
$\iota(f), \iota(g) : \STop \to \STop'$ of the possible-worlds
topoi.  Since distinct model maps induce distinct inverse image
functors on the sheaf categories (they disagree on at least one
stalk), $\iota(f) \neq \iota(g)$.

\emph{Preservation of truth.}  In Montague semantics,
$\llbracket\varphi\rrbracket(w) = \mathrm{true}$ iff $\varphi$ is
true at world~$w$.  In the JuGeo embedding, this corresponds to the
global section $s(w) = \llbracket\varphi\rrbracket(w)$ being
non-empty, i.e., $\Gamma(\STop, \mathcal{F}_\varphi) \neq \emptyset$.
Since $H^1 = 0$ (trivial cohomology), the fundamental theorem gives
$P \models \varphi$ iff $H^0 \neq 0$, which is exactly the Montague
truth condition.
\end{proof}

\begin{example}
The sentence ``Every man walks'' has Montague translation
$\forall x.\, \mathrm{man}(x) \to \mathrm{walk}(x)$ of type~$t$.
In the embedding: $\STop = \Sh(W)$, $A = \Pi(x:e).\,
\mathrm{man}(x) \to \mathrm{walk}(x)$, $T = \gone$ if the
sentence is true at all worlds, and $H^1 = 0$ (no cohomological
obstruction).
\end{example}

%% ── §19.2  DRT/SDRT ────────────────────────────────────────────────
\section{JuGeo subsumes DRT/SDRT}
\label{sec:sub-drt}

\begin{definition}[DRT]
\label{def:drt}
A \emph{Discourse Representation Structure} (DRS) is a pair
$K = (\mathcal{R}, \mathcal{C})$ where $\mathcal{R}$ is a set of
discourse referents and $\mathcal{C}$ is a set of conditions.
A DRS is \emph{satisfied} by an embedding $f : \mathcal{R} \to D$
into a model domain $D$ if all conditions in $\mathcal{C}$ are true
under~$f$.  SDRT extends DRT with rhetorical relations between
discourse segments.
\end{definition}

\begin{construction}[Embedding $\iota_{\mathrm{DRT}} :
  \mathbf{DRT} \to \mathbf{Judg}_4$]
\label{con:drt-embed}
Given a DRS $K = (\mathcal{R}, \mathcal{C})$:
\begin{itemize}[nosep]
\item $\STop$: The site $\mathcal{C}_{\mathrm{disc}}$ has objects =
  discourse segments $d_1, \ldots, d_n$ and morphisms = anaphoric
  accessibility relations.  The topology is generated by covers
  $\{d_i \to d_{\mathrm{all}}\}$.  The presheaf
  $\mathcal{F}(d_k) = \{f : \mathcal{R}_k \to D\}$ assigns
  possible embeddings.

\item $\PStr$: Define a \emph{discourse monad} type constructor
  $\mathrm{DMon}(A)$ that encapsulates a set of discourse referents
  and a value of type $A$ depending on them.  A DRS becomes a
  term $E_K : \mathrm{DMon}(\mathrm{Prop})$ where
  $\mathrm{Prop} = \Sigma(c : \mathcal{C}).\, \mathrm{Sat}(c, f)$.
  Context: $\Gamma = (D : \Type,\; f : \mathcal{R} \to D)$.

\item $\CSpec$: The \v{C}ech complex is computed on the discourse
  site.  $H^1$ vanishes iff anaphoric references resolve consistently
  across all discourse segments.

\item $\GEval$: $G = \mathbf{2}$.  Grade $\gone$ iff all DRS
  conditions are satisfied.
\end{itemize}
For SDRT: rhetorical relations $R(d_i, d_j)$ become morphisms in the
site category, with the topology encoding which rhetorical structures
constitute complete discourse covers.
\end{construction}

\begin{theorem}
\label{thm:drt}
The embedding $\iota_{\mathrm{DRT}}$ is faithful and preserves
DRS satisfaction.
\end{theorem}

\begin{proof}
\emph{Faithfulness.}  Distinct DRS morphisms (embeddings $f \neq g$)
yield distinct context morphisms in $\PStr$ (the substitutions
$\sigma_f \neq \sigma_g$ disagree on at least one referent),
hence distinct $\mathbf{Judg}_4$ morphisms.

\emph{Preservation of satisfaction.}  A DRS $K$ is satisfied by
$f$ iff all conditions hold under~$f$.  In the embedding, this
corresponds to the existence of a global section of $\mathcal{F}$
over the discourse site---i.e., a compatible family of local
embeddings that agree on shared referents.  The fundamental
theorem identifies this with $H^1 = 0$.  Since DRT
satisfaction is precisely the existence of a consistent global
embedding, the two notions coincide.
\end{proof}

\begin{example}
``A man$_x$ entered.  He$_x$ smiled.''  The discourse site has two
segments $d_1, d_2$ with accessibility $d_1 \to d_2$.  The DRS has
$\mathcal{R} = \{x\}$, $\mathcal{C} = \{\mathrm{man}(x),
\mathrm{enter}(x), \mathrm{smile}(x)\}$.  Local sections assign
embeddings $f_1 : \{x\} \to D$ and $f_2 : \{x\} \to D$.
Compatibility on the overlap $d_1 \cap d_2$ requires $f_1(x) =
f_2(x)$, i.e., the pronoun ``he'' corefers with ``a man.''
\end{example}

%% ── §19.3  MLTT ────────────────────────────────────────────────────
\section{JuGeo subsumes MLTT}
\label{sec:sub-mltt}

\begin{definition}[MLTT]
\label{def:mltt}
Martin-L\"of Type Theory (MLTT) consists of:
\begin{enumerate}[label=(\roman*),nosep]
\item Dependent function types $\Pi(x:A).\,B(x)$;
\item Dependent pair types $\Sigma(x:A).\,B(x)$;
\item Identity types $\mathrm{Id}_A(a,b)$;
\item A universe hierarchy $\Type_0 : \Type_1 : \Type_2 : \cdots$;
\item Inductive types (including $\mathbb{N}$, $W$-types).
\end{enumerate}
Judgments have the form $\Gamma \vdash e : A$.
\end{definition}

\begin{construction}[Embedding $\iota_{\mathrm{MLTT}} :
  \mathbf{MLTT} \to \mathbf{Judg}_4$]
\label{con:mltt-embed}
Given an MLTT judgment $\Gamma \vdash e : A$:
\begin{itemize}[nosep]
\item $\STop$: The site is the terminal category $\mathbf{1}$ with
  the trivial (maximal) topology.  The topos
  $\Sh(\mathbf{1}) \cong \Set$.  The presheaf
  $\mathcal{F}(*) = \llbracket A \rrbracket$ is the set-theoretic
  interpretation of~$A$.

\item $\PStr$: Exactly the MLTT judgment $\Gamma \vdash e : A$ with
  its full derivation tree, type formers, and reduction rules.

\item $\CSpec$: The cover is the trivial cover $\{* \to *\}$.
  The \v{C}ech complex is $C^0 = \mathcal{F}(*)$,
  $C^n = 0$ for $n \geq 1$.  All higher cohomology vanishes:
  $H^n = 0$ for $n \geq 1$.

\item $\GEval$: $G = \mathbf{2}$.  Grade $\gone$ iff
  the judgment is derivable, $\gzero$ otherwise.
\end{itemize}
\end{construction}

\begin{theorem}
\label{thm:mltt}
The embedding $\iota_{\mathrm{MLTT}}$ is faithful, full on the binary
fragment, and preserves derivability.
\end{theorem}

\begin{proof}
\emph{Faithfulness.}  MLTT morphisms are substitutions
$\sigma : \Gamma \to \Gamma'$ with $A'[\sigma] \equiv A$ and
$e'[\sigma] \equiv e$.  These are exactly the context morphisms
in $\PStr$, hence distinct substitutions give distinct
$\mathbf{Judg}_4$ morphisms.

\emph{Fullness on the binary fragment.}  Every $\mathbf{Judg}_4$
morphism in the image of $\iota_{\mathrm{MLTT}}$ has trivial topos
component (since $\STop = \Set$), trivial cohomology, and binary
grade.  The only non-trivial data is the context morphism, which is
an MLTT substitution.  Hence every morphism in the binary-graded,
flat-site fragment of $\mathbf{Judg}_4$ comes from an MLTT morphism.

\emph{Preservation of derivability.}  $\Gamma \vdash e : A$ is
derivable in MLTT iff the embedded judgment has grade $\gone$, since
the grade is $\gone$ exactly when the derivation tree exists.
The cohomology vanishes trivially ($H^1 = 0$ on the terminal site),
so the fundamental theorem reduces to: a global section exists iff
a term of type $A$ exists in context $\Gamma$.
\end{proof}

%% ── §19.4  OT ──────────────────────────────────────────────────────
\section{JuGeo subsumes Optimality Theory}
\label{sec:sub-ot}

\begin{definition}[OT]
\label{def:ot}
An Optimality Theory grammar is a triple
$(\mathrm{Gen}, \mathrm{Con}, \gg)$ where $\mathrm{Gen}$ is a
generator function mapping inputs to candidate sets,
$\mathrm{Con} = \{C_1, \ldots, C_k\}$ is a finite set of ranked
constraints with $C_1 \gg C_2 \gg \cdots \gg C_k$, and the
\emph{optimal candidate} is
$\arg\min_{\mathrm{cand}} (C_1(\mathrm{cand}), \ldots,
C_k(\mathrm{cand}))$ under lexicographic ordering.
\end{definition}

\begin{construction}[Embedding $\iota_{\mathrm{OT}} :
  \mathbf{OT} \to \mathbf{Judg}_4$]
\label{con:ot-embed}
Given an OT grammar $(\mathrm{Gen}, \mathrm{Con}, \gg)$ and an
input~$i$:
\begin{itemize}[nosep]
\item $\STop$: The site has a single stratum (the phonological
  level).  Objects are candidates
  $\mathrm{Gen}(i) = \{o_1, \ldots, o_m\}$.
  Morphisms are identity only.  Topology: trivial.

\item $\PStr$: Context $\Gamma = (i : \mathrm{Input})$.
  For each candidate $o_j$, the type
  $A_j = \Pi(k : [K]).\, C_k(o_j) \leq C_k(o_l)$ (for all
  competitors $o_l$) encodes optimality.  The evidence term
  witnesses the constraint violations.

\item $\CSpec$: Trivial (single stratum, no non-trivial covers).

\item $\GEval$: The grade semiring is $G = \mathbb{N}^k$ with
  lexicographic order.  Define $\otG$ as componentwise addition and
  $\opG$ as componentwise minimum.  The grade of candidate $o_j$ is
  $\mu(o_j) = (C_1(o_j), \ldots, C_k(o_j))$.  The optimal
  candidate has the maximal grade under the reversed lexicographic
  order (fewest violations).

  Equivalently, assign exponential weights $w_k = N^{k-1}$ for
  large $N$ and compute
  $\mu(o_j) = \sum_k w_k \cdot C_k(o_j)$,
  reducing to a single-valued grade with lexicographic
  tie-breaking.
\end{itemize}
\end{construction}

\begin{theorem}
\label{thm:ot}
The embedding $\iota_{\mathrm{OT}}$ is faithful and the optimal
candidate in OT corresponds to the judgment with maximal harmony.
\end{theorem}

\begin{proof}
\emph{Faithfulness.}  Distinct OT grammars have distinct constraint
rankings, which yield distinct grade semirings (different orderings
on $\mathbb{N}^k$), hence distinct $\GEval$ components.

\emph{Optimality $=$ harmony.}  The OT optimal candidate minimizes
violation counts lexicographically.  In the JuGeo embedding, the
harmony of candidate~$o_j$ is
\[
  h(o_j) = \bigotimes_{k=1}^{K} w_k \cdot g_k(o_j),
\]
where $g_k(o_j) = C_k(o_j)$ is the constraint violation count.
The candidate with minimal violation vector has maximal harmony
(since $\otG$ is the bottleneck operation and we use the reversed
ordering).  Thus OT evaluation is a special case of $\GEval$
computation.
\end{proof}

\begin{example}
Input /pat/ with constraints
$\mathrm{Onset} \gg \mathrm{NoCoda} \gg \mathrm{Faith}$.
Candidates: [pat], [pa.ta], [pa].
Grade of [pat]: $(0, 1, 0)$.  Grade of [pa.ta]: $(0, 0, 1)$.
Optimal: [pa.ta] (fewer high-ranked violations).
In JuGeo: harmony $h$ is maximal for [pa.ta].
\end{example}

%% ── §19.5  Linear Logic ────────────────────────────────────────────
\section{JuGeo subsumes Linear Logic}
\label{sec:sub-linear}

\begin{definition}[Linear Logic]
\label{def:linear}
Girard's Linear Logic (LL) is a substructural logic with:
\begin{enumerate}[label=(\roman*),nosep]
\item Multiplicative connectives: $\otimes$ (tensor),
  $\multimap$ (linear implication), $\mathbf{1}$ (unit);
\item Additive connectives: $\oplus$ (plus), $\&$ (with),
  $\mathbf{0}$ (zero), $\top$;
\item Exponentials: $!A$ (of course), $?A$ (why not);
\item No unrestricted weakening or contraction---resources are
  tracked exactly.
\end{enumerate}
\end{definition}

\begin{construction}[Embedding $\iota_{\mathrm{LL}} :
  \mathbf{LL} \to \mathbf{Judg}_4$]
\label{con:linear-embed}
Given an LL sequent $\Gamma \vdash A$:
\begin{itemize}[nosep]
\item $\STop$: Terminal site $\mathbf{1}$ with $\Sh(\mathbf{1})
  \cong \Set$.

\item $\PStr$: The LL proof is embedded into dependent type theory
  via the standard Curry--Howard correspondence for linear logic:
  $A \otimes B$ becomes $\Sigma_{\mathrm{lin}}(x:A).\,B$,
  $A \multimap B$ becomes $\Pi_{\mathrm{lin}}(x:A).\,B$,
  and the exponential $!A$ becomes the unrestricted modality.
  The context~$\Gamma$ tracks linear resources.

\item $\CSpec$: Trivial (flat site).

\item $\GEval$: The grade semiring captures resource usage.
  Set $G = \mathbb{N}$ (counting resource uses).
  $\otG = +$ (combining resources), $\opG = \max$ (alternative
  resource usage), $\gone = 0$ (no resources used),
  $\gzero = \infty$ (impossible).
  The LL constraint ``use each resource exactly once'' becomes
  the grade condition $\mu(\mathcal{D}) = 1$ for each linear
  hypothesis.

  Alternatively, for the full multiplicative-additive structure:
  $G = (\mathbb{N}^n, +, \max, \vec{0}, \vec{\infty})$ where $n$
  counts the number of distinct resources, and the grade tracks
  the exact usage vector.
\end{itemize}
\end{construction}

\begin{theorem}
\label{thm:linear}
The embedding $\iota_{\mathrm{LL}}$ is faithful and preserves
provability.
\end{theorem}

\begin{proof}
\emph{Faithfulness.}  LL proofs are derivation trees with strict
resource discipline.  Distinct proofs have distinct resource usage
patterns, hence distinct grade evaluations.  Since the $\PStr$
component records the full derivation and the $\GEval$ component
records the resource usage, distinct LL morphisms yield distinct
$\mathbf{Judg}_4$ morphisms.

\emph{Preservation of provability.}  An LL sequent
$\Gamma \vdash A$ is provable iff there exists a derivation tree
using each linear hypothesis exactly once.  In the embedding, this
corresponds to the existence of a term $E$ of type $A$ in the
linear context $\Gamma$ with grade $\mu(\mathcal{D}) = \vec{1}$
(each resource used exactly once).  The grade condition
$\mu(\mathcal{D}) = \vec{1}$ is equivalent to the linear logic
resource discipline.
\end{proof}

%% ── §19.6  Generative Lexicon ──────────────────────────────────────
\section{JuGeo subsumes the Generative Lexicon}
\label{sec:sub-gl}

\begin{definition}[Generative Lexicon]
\label{def:gl}
Pustejovsky's Generative Lexicon (GL) assigns to each lexical item
a \emph{qualia structure} $\mathrm{Qualia}(w) = (Q_F, Q_C, Q_T,
Q_A)$ consisting of:
\begin{enumerate}[label=(\roman*),nosep]
\item $Q_F$: formal quale (the type/category);
\item $Q_C$: constitutive quale (parts/material);
\item $Q_T$: telic quale (purpose/function);
\item $Q_A$: agentive quale (origin/creation).
\end{enumerate}
\emph{Type coercion} occurs when a verb selects an argument via
one quale (e.g., ``begin the book'' coerces ``book'' via $Q_T$
to ``reading the book'').
\end{definition}

\begin{construction}[Embedding $\iota_{\mathrm{GL}} :
  \mathbf{GL} \to \mathbf{Judg}_4$]
\label{con:gl-embed}
Given a lexical item $w$ with qualia $\mathrm{Qualia}(w)$:
\begin{itemize}[nosep]
\item $\STop$: The site is the lexical category site, with objects
  = lexical items and morphisms = sense relations (hypernymy,
  meronymy, etc.).

\item $\PStr$: Define the \emph{qualia type constructor}
  \[
    \mathrm{Qualia}(C) = \Sigma(q_F : \mathrm{Form}(C)).\,
    \Sigma(q_C : \mathrm{Const}(C)).\,
    \Sigma(q_T : \mathrm{Telic}(C)).\,
    \mathrm{Agent}(C)
  \]
  as a dependent record type.  Coercion is a function
  $\mathrm{coerce}_T : \mathrm{Qualia}(C) \to C_T$ that
  projects the telic quale and interprets it as a type.

\item $\CSpec$: Trivial for single-word analysis.  For
  compositional semantics across a sentence, the cohomology
  detects coercion conflicts (incompatible qualia across
  composition).

\item $\GEval$: $G = \mathbf{2}$.  Grade $\gone$ if the qualia
  structure is well-formed and coercion succeeds.
\end{itemize}
\end{construction}

\begin{theorem}
\label{thm:gl}
The embedding $\iota_{\mathrm{GL}}$ is faithful and preserves
type coercion.
\end{theorem}

\begin{proof}
\emph{Faithfulness.}  Distinct GL lexical entries have distinct
qualia structures, which yield distinct $\Sigma$-types in $\PStr$.
Distinct coercion operations yield distinct projection functions.

\emph{Preservation of coercion.}  GL coercion ``begin the book''
$\rightsquigarrow$ ``begin reading the book'' is modeled by
selecting $Q_T(\text{book}) = \text{read}$ from the qualia structure.
In the JuGeo embedding, this is the projection
$\pi_T : \mathrm{Qualia}(\text{book}) \to \mathrm{Telic}(\text{book})$
applied to the term $E_{\text{book}} : \mathrm{Qualia}(\text{book})$.
The coerced type $\mathrm{Telic}(\text{book}) = \text{read}(\text{book})$
is the telic quale, and the derivation witnessing
$\Gamma \vdash \mathrm{begin}(\pi_T(E_{\text{book}})) :
\mathrm{Event}$ is exactly the GL coercion mechanism.
\end{proof}

%% ── §19.7  HoTT ───────────────────────────────────────────────────
\section{JuGeo subsumes HoTT}
\label{sec:sub-hott}

\begin{definition}[HoTT identity types]
\label{def:hott}
In Homotopy Type Theory, the identity type $\mathrm{Id}_A(a,b)$ for
$a, b : A$ satisfies:
\begin{enumerate}[label=(\roman*),nosep]
\item Reflexivity: $\mathrm{refl}_a : \mathrm{Id}_A(a,a)$;
\item Symmetry: $p : \mathrm{Id}_A(a,b) \vdash p^{-1} :
  \mathrm{Id}_A(b,a)$;
\item Transitivity: $p : \mathrm{Id}_A(a,b),\; q :
  \mathrm{Id}_A(b,c) \vdash p \cdot q : \mathrm{Id}_A(a,c)$;
\item These form an $\infty$-groupoid structure on identity proofs.
\end{enumerate}
\end{definition}

\begin{construction}[Embedding $\iota_{\mathrm{HoTT}} :
  \mathbf{HoTT} \to \mathbf{Judg}_4$]
\label{con:hott-embed}
Given a HoTT judgment $\Gamma \vdash e : A$:
\begin{itemize}[nosep]
\item $\STop$, $\CSpec$: As for MLTT (\cref{con:mltt-embed}):
  terminal site, trivial cohomology.

\item $\PStr$: The full HoTT type theory, including higher inductive
  types, univalence, and the $\infty$-groupoid structure.

\item $\GEval$: The graded identity type $\GId_g(A,a,b)$ at grade
  $g = \gone$ is set to be exactly the HoTT identity type:
  \[
    \GId_{\gone}(A,a,b) \;\coloneqq\; \mathrm{Id}_A(a,b).
  \]
  At sub-maximal grades $g < \gone$, $\GId_g(A,a,b)$ is a
  relaxed ``near-identity'' type (inhabited when $a$ and $b$ are
  $g$-close but not necessarily equal).  The grade semiring is
  $G = [0,1]$ with $\otG = \min$, $\opG = \max$.
\end{itemize}
\end{construction}

\begin{theorem}
\label{thm:hott}
The embedding $\iota_{\mathrm{HoTT}}$ is faithful and
$\GId_{\gone}(A,a,b) = \mathrm{Id}_A(a,b)$.
\end{theorem}

\begin{proof}
\emph{Faithfulness.}  Follows from faithfulness of the MLTT embedding
(\cref{thm:mltt}), since HoTT extends MLTT.

\emph{Identity type recovery.}
By definition, $\GId_{\gone}(A,a,b) = \mathrm{Id}_A(a,b)$.  We
verify that the graded groupoid structure at grade $\gone$
coincides with the HoTT groupoid structure:
\begin{enumerate}[label=(\roman*),nosep]
\item Reflexivity: $\mathrm{refl}_a : \GId_{\gone}(A,a,a)$ is
  exactly $\mathrm{refl}_a : \mathrm{Id}_A(a,a)$.
\item Symmetry: the graded inverse at $g = \gone$ is HoTT path
  inversion.
\item Transitivity: the graded composition at $g = \gone$ is HoTT
  path concatenation.
\item Higher structure: the graded groupoid laws at $g = \gone$
  are the HoTT groupoid laws (associativity, left/right unit,
  inverse laws hold up to higher identity).
\end{enumerate}
Thus the HoTT $\infty$-groupoid structure is the $g = \gone$ slice
of the graded groupoid.
\end{proof}

\begin{remark}
At grades $g < \gone$, $\GId_g$ provides structure not present in
HoTT: approximate or graded identifications.  This is what makes
the JuGeo framework a \emph{proper} extension (\cref{thm:proper}).
\end{remark}

%% ── §19.8  Prior's Tense Logic ─────────────────────────────────────
\section{JuGeo subsumes Prior's Tense Logic}
\label{sec:sub-prior}

\begin{definition}[Tense Logic]
\label{def:prior}
Prior's basic tense logic extends propositional logic with two
modal operators:
\begin{enumerate}[label=(\roman*),nosep]
\item $P\varphi$ (``it was the case that $\varphi$'');
\item $F\varphi$ (``it will be the case that $\varphi$'').
\end{enumerate}
Models are triples $(T, <, V)$ where $(T, <)$ is a strict linear
order (time) and $V : \mathrm{Prop} \to \mathcal{P}(T)$ is a
valuation.
\end{definition}

\begin{construction}[Embedding $\iota_{\mathrm{TL}} :
  \mathbf{TL} \to \mathbf{Judg}_4$]
\label{con:prior-embed}
Given a tense-logical model $(T, <, V)$ and a formula $\varphi$:
\begin{itemize}[nosep]
\item $\STop$: The site $\mathcal{C}_T$ has objects = intervals
  $[a,b] \subseteq T$ and morphisms = inclusions.  The topology is
  generated by covers $\{[a,c], [c,b]\} \to [a,b]$ (interval
  bisections).  The presheaf $\mathcal{F}([a,b]) = \{t \in [a,b]
  \mid V(\varphi, t) = \mathrm{true}\}$.

\item $\PStr$: The temporal operators become type constructors:
  $P(A) = \Sigma(t' : T).\, (t' < t) \times A(t')$ and
  $F(A) = \Sigma(t' : T).\, (t < t') \times A(t')$.  The
  context is $\Gamma = (t : T)$.

\item $\CSpec$: The \v{C}ech complex on the temporal site detects
  whether a temporally local property holds globally over a time
  interval.

\item $\GEval$: $G = \mathbf{2}$.  Grade $\gone$ iff $\varphi$ is
  true at the evaluation time.
\end{itemize}
\end{construction}

\begin{theorem}
\label{thm:prior}
The embedding $\iota_{\mathrm{TL}}$ is faithful and preserves
tense-logical validity.
\end{theorem}

\begin{proof}
\emph{Faithfulness.}  Distinct tense models $(T, <, V) \neq
(T', <', V')$ yield distinct sites (different interval categories or
different valuations), hence distinct $\STop$ components.  Distinct
tense formulas yield distinct types in $\PStr$.

\emph{Preservation of validity.}  A tense formula $\varphi$ is valid
at time $t$ in model $(T, <, V)$ iff $V(\varphi, t) = \mathrm{true}$.
In the embedding, this is the existence of a global section of
$\mathcal{F}$ over the point $\{t\}$ (which is trivially a cover of
itself), i.e., $\mathcal{F}(\{t\}) \neq \emptyset$.

For temporally extended validity (``$\varphi$ holds throughout
$[a,b]$''), the fundamental theorem applies: $\varphi$ holds on all
of $[a,b]$ iff $H^1 = 0$ on the interval cover, iff local truth at
sub-intervals glues to global truth.

The tense operators are preserved by construction:
$P\varphi$ is true at $t$ iff
$\exists t' < t.\, V(\varphi, t')$ iff the $\Sigma$-type
$\Sigma(t':T).\,(t' < t) \times \mathcal{F}(\{t'\})$ is inhabited.
\end{proof}

%% ── §19.9  Proper Extension ────────────────────────────────────────
\section{JuGeo is a proper extension}
\label{sec:proper}

\begin{theorem}[Proper extension]
\label{thm:proper}
$\mathbf{Judg}_4$ is not equivalent to any of the subsumed
frameworks.  That is, the embeddings $\iota_{\mathrm{MG}},
\iota_{\mathrm{DRT}}, \iota_{\mathrm{MLTT}}, \iota_{\mathrm{OT}},
\iota_{\mathrm{LL}}, \iota_{\mathrm{GL}}, \iota_{\mathrm{HoTT}},
\iota_{\mathrm{TL}}$ are all proper (faithful but not essentially
surjective).
\end{theorem}

\begin{proof}
We exhibit a single judgment $\JJ^* \in \mathbf{Judg}_4$ that lies
outside the image of every embedding.  Define $\JJ^*$ as follows:
\begin{itemize}[nosep]
\item $\STop^*$: a non-trivial site with three strata
  (phonological, semantic, pragmatic) and covers linking all three;
\item $\PStr^*$: a graded judgment
  $\Gamma \vdash_{0.7} E : \mathrm{Poem}$ with fractional grade;
\item $\CSpec^*$: non-trivial $H^1 \neq 0$ (the poem has a forced
  rhyme creating a cohomological obstruction);
\item $\GEval^*$: the grade semiring $[0,1]$ with the bottleneck
  product, evaluating to $T = 0.7$ (the naturalness stratum is the
  bottleneck).
\end{itemize}

This judgment cannot lie in the image of:
\begin{enumerate}[label=(\roman*),nosep]
\item $\iota_{\mathrm{MG}}$: Montague has binary grades and no
  multi-stratal structure ($|G| = 2$, but $\JJ^*$ has $T = 0.7$).
\item $\iota_{\mathrm{DRT}}$: DRT has binary grades and no
  cohomological structure for poetry.
\item $\iota_{\mathrm{MLTT}}$: MLTT has trivial site and binary
  grades.
\item $\iota_{\mathrm{OT}}$: OT has a single stratum and integer
  violation counts (not continuous grades).
\item $\iota_{\mathrm{LL}}$: Linear Logic has resource-counting
  grades ($\mathbb{N}$), not continuous $[0,1]$, and trivial site.
\item $\iota_{\mathrm{GL}}$: GL has binary grades and no
  cohomological structure.
\item $\iota_{\mathrm{HoTT}}$: HoTT has trivial site and no
  multi-stratal harmony.
\item $\iota_{\mathrm{TL}}$: Tense Logic has binary grades and
  temporal site only (not multi-stratal).
\end{enumerate}
Since $\JJ^*$ uses continuous grades, non-trivial cohomology,
\emph{and} multi-stratal structure simultaneously, no single
subsumed framework can express it.
\end{proof}

%%%%%%%%%%%%%%%%%%%%%%%%%%%%%%%%%%%%%%%%%%%%%%%%%%%%%%%%%%%%%%%%%%%%%%%%
%%  Chapter 20 — Descent and Harmony Proofs
%%%%%%%%%%%%%%%%%%%%%%%%%%%%%%%%%%%%%%%%%%%%%%%%%%%%%%%%%%%%%%%%%%%%%%%%
\chapter{Descent and Harmony Proofs}
\label{ch:descent-proofs}

This chapter contains the full proofs of the descent theorem, the
harmony existence and uniqueness theorems, the bottleneck principle,
the $\FMLM$ convergence theorem, and the descent-as-optimisation
correspondence.

%% ── §20.1  The Descent Theorem ──────────────────────────────────────
\section{The descent theorem}
\label{sec:descent-thm}

We first establish the fundamental theorem of Judgment Geometry: the
equivalence between global validity and the vanishing of the first
\v{C}ech cohomology.

\begin{definition}[\v{C}ech complex]
\label{def:cech-complex}
Let $(\mathcal{C}, \mathcal{J})$ be a site, $\mathcal{F}$ a presheaf
on $\mathcal{C}$, and $\mathcal{U} = \{U_i \to X\}_{i \in I}$ a
covering family of an object $X$.  The \v{C}ech complex
$\Cech^\bullet(\mathcal{U}, \mathcal{F})$ is:
\[
  \Cech^n(\mathcal{U}, \mathcal{F})
  = \prod_{i_0 < \cdots < i_n}
    \mathcal{F}(U_{i_0} \times_X \cdots \times_X U_{i_n})
\]
with coboundary maps $\delta^n : \Cech^n \to \Cech^{n+1}$ given by
\[
  (\delta^n s)_{i_0 \ldots i_{n+1}}
  = \sum_{k=0}^{n+1} (-1)^k\, s_{i_0 \ldots \hat{i}_k \ldots i_{n+1}}
    \big|_{U_{i_0} \times_X \cdots \times_X U_{i_{n+1}}}.
\]
The cohomology groups are $\check{H}^n(\mathcal{U}, \mathcal{F})
= \ker \delta^n / \mathrm{im}\, \delta^{n-1}$.
\end{definition}

\begin{definition}[Descent datum]
\label{def:descent-datum}
A \emph{descent datum} for $\mathcal{F}$ on the cover $\mathcal{U}$
is a family of local sections $\{s_i \in \mathcal{F}(U_i)\}_{i \in I}$
satisfying the \emph{cocycle condition}: for all $i, j \in I$,
\[
  s_i|_{U_i \times_X U_j} = s_j|_{U_i \times_X U_j}.
\]
Equivalently, a descent datum is a 0-cocycle $s \in \ker \delta^0$.
\end{definition}

\begin{theorem}[Fundamental Theorem of Judgment Geometry]
\label{thm:fundamental}
Let $(\mathcal{C}, \mathcal{J})$ be a site, $\mathcal{F}$ a presheaf,
$\mathcal{U}$ a cover of $X$, and $\varphi$ a property encoded by
$\mathcal{F}$.  Then:
\[
  X \models \varphi
  \quad\Longleftrightarrow\quad
  \check{H}^1(\mathcal{U}, \mathcal{F}) = 0
  \;\;\text{and}\;\;
  \check{H}^0(\mathcal{U}, \mathcal{F}) \neq 0.
\]
In particular, if $\mathcal{F}$ is a sheaf, then $X \models \varphi$
if and only if a compatible family of local sections (each witnessing
$\varphi$ on its region) exists.
\end{theorem}

\begin{proof}
$(\Rightarrow)$\; Suppose $X \models \varphi$, i.e., there exists a
global section $s \in \mathcal{F}(X)$ witnessing~$\varphi$.
Restricting $s$ to each $U_i$ gives $s_i = s|_{U_i}$.  This family
is automatically compatible:
$s_i|_{U_{ij}} = s|_{U_{ij}} = s_j|_{U_{ij}}$, so $(s_i)$ is a
0-cocycle and $\check{H}^0 \neq 0$.

For $\check{H}^1 = 0$: let $z = (z_{ij}) \in \ker \delta^1$ be any
1-cocycle.  We must show $z$ is a coboundary.  The cocycle condition
gives $z_{ij} + z_{jk} = z_{ik}$ on triple overlaps.  Define
$t_i = z_{i0}$ for a fixed index~$0$.  Then
\[
  (\delta^0 t)_{ij} = t_j - t_i = z_{j0} - z_{i0} = z_{ij}
\]
by the cocycle condition $z_{i0} + z_{0j} = z_{ij}$ (using
$z_{0j} = -z_{j0}$ by antisymmetry).  Hence $z = \delta^0 t$, so
$\check{H}^1 = 0$.

$(\Leftarrow)$\;
Suppose $\check{H}^0 \neq 0$ and $\check{H}^1 = 0$.  The condition
$\check{H}^0 \neq 0$ gives a compatible family $(s_i)_{i \in I}$
with $s_i|_{U_{ij}} = s_j|_{U_{ij}}$ for all $i, j$.  If
$\mathcal{F}$ is a sheaf, the sheaf condition states precisely that
such a compatible family glues to a unique global section
$s \in \mathcal{F}(X)$.  Hence $X \models \varphi$.

If $\mathcal{F}$ is only a presheaf, the vanishing $\check{H}^1 = 0$
is an additional hypothesis ensuring that the obstruction to gluing
(which lives in $\check{H}^1$) vanishes.  Concretely: suppose we have
local sections $s_i$ that are \emph{almost} compatible, differing by
a discrepancy $d_{ij} = s_i|_{U_{ij}} - s_j|_{U_{ij}}$.  The
discrepancy $(d_{ij})$ is a 1-cocycle (it satisfies the cocycle
condition on triple overlaps by a telescoping argument).  If
$\check{H}^1 = 0$, then $(d_{ij}) = (\delta^0 t)_{ij}$ for some
$(t_i)$, and the corrected sections $s_i' = s_i - t_i$ are
compatible: $s_i'|_{U_{ij}} - s_j'|_{U_{ij}} = d_{ij} -
(t_i|_{U_{ij}} - t_j|_{U_{ij}}) = d_{ij} - d_{ij} = 0$.
The corrected family glues (by the sheaf condition applied to the
sheafification of $\mathcal{F}$).
\end{proof}

\begin{corollary}[Verification $= H^0$]
\label{cor:verification}
$X \models \varphi$ if and only if
$\Gamma(X, \mathcal{F}_\varphi) \neq \emptyset$
(the global sections functor is non-empty).  For an acyclic cover
(where \v{C}ech and derived cohomology agree), this is identified
with $\check{H}^0(\mathcal{U}, \mathcal{F}) \neq \emptyset$.
\end{corollary}

\begin{proof}
The global sections functor $\Gamma(X, -)$ is
$\Hom_{\PSh(\mathcal{C})}(\mathbf{y}(X), -)$.  A global section of
$\mathcal{F}_\varphi$ is precisely an element of $\mathcal{F}(X)$
witnessing~$\varphi$.  By \cref{thm:fundamental}, this exists iff
$H^0 \neq 0$ and $H^1 = 0$.  For acyclic covers, the Leray spectral
sequence degenerates: $\check{H}^n(\mathcal{U}, \mathcal{F}) \cong
H^n(X, \mathcal{F})$ for all $n$, so
$\Gamma(X, \mathcal{F}) \cong \check{H}^0$.
\end{proof}

\subsection{The \v{C}ech-to-derived spectral sequence}

\begin{proposition}[\v{C}ech-to-derived comparison]
\label{prop:cech-derived}
For any cover $\mathcal{U}$ of $X$ and any sheaf $\mathcal{F}$, there
is a spectral sequence
\[
  E_2^{p,q} = \check{H}^p(\mathcal{U}, \underline{H}^q(\mathcal{F}))
  \;\Longrightarrow\;
  H^{p+q}(X, \mathcal{F}),
\]
where $\underline{H}^q(\mathcal{F})$ is the presheaf
$U \mapsto H^q(U, \mathcal{F}|_U)$.  When the cover is acyclic
($H^q(U_{i_0 \cdots i_p}, \mathcal{F}) = 0$ for $q > 0$ and all
multi-intersections), this degenerates at $E_2$ and gives
$\check{H}^n(\mathcal{U}, \mathcal{F}) \cong H^n(X, \mathcal{F})$.
\end{proposition}

\begin{proof}
The double complex $C^{p,q} = \Cech^p(\mathcal{U}, \mathcal{I}^q)$,
where $\mathcal{F} \to \mathcal{I}^\bullet$ is an injective
resolution, has two filtrations.  Filtering by rows gives
$E_1^{p,q} = \Cech^p(\mathcal{U}, \underline{H}^q)$, whose
cohomology is $E_2^{p,q} = \check{H}^p(\mathcal{U},
\underline{H}^q)$.  Filtering by columns gives $E_1^{p,q} =
H^q(\prod U_{i_0 \cdots i_p}, \mathcal{F})$; when the cover is
acyclic, $E_1^{p,q} = 0$ for $q > 0$, so the column filtration
degenerates at $E_1$ and both filtrations converge to
$H^{p+q}(X, \mathcal{F})$.  The isomorphism
$\check{H}^n \cong H^n$ follows \cite{weibel1994}.
\end{proof}

%% ── §20.2  Harmony Existence ────────────────────────────────────────
\section{Harmony existence and uniqueness}
\label{sec:harmony}

\begin{definition}[Harmony function]
\label{def:harmony}
Let $\mathcal{T}_1, \ldots, \mathcal{T}_N$ be theories (theory slots)
with stratal grades $g_1, \ldots, g_N$ and weights
$w_1, \ldots, w_N \in \Grade$.  The \emph{harmony function}
$h : \mathrm{Sec}(\mathcal{F}) \to \Grade$ is defined by
\[
  h(\sigma) = \bigotimes_{k=1}^{N} w_k \otG g_k(\sigma),
\]
where $g_k(\sigma)$ is the grade of $\sigma$ with respect to
theory~$\mathcal{T}_k$, and $\mathrm{Sec}(\mathcal{F})$ is the set
of (candidate) sections of the judgment presheaf.
\end{definition}

\begin{definition}[Harmony lattice]
\label{def:harmony-lattice}
The \emph{harmony lattice} is the complete lattice
$(\mathrm{Sec}(\mathcal{F}), \sqsubseteq)$ where
$\sigma \sqsubseteq \sigma'$ iff $h(\sigma) \leq h(\sigma')$
(ordered by harmony value).
\end{definition}

\begin{theorem}[Harmony fixpoint existence]
\label{thm:harmony-fix}
Let $h : \mathrm{Sec}(\mathcal{F}) \to \Grade$ be the harmony
function.  Suppose:
\begin{enumerate}[label=(\roman*),nosep]
\item The number of strata is finite: $N < \infty$.
\item The evaluation $g_k$ is continuous (Scott-continuous) for each
  theory $\mathcal{T}_k$.
\item The grade semiring $(\Grade, \leq)$ is a complete lattice.
\end{enumerate}
Then the refinement operator
$R(\sigma) = \arg\max_{\sigma'} h(\sigma')$ subject to
$\sigma' \sqsubseteq \sigma$ has a fixpoint.
\end{theorem}

\begin{proof}
We apply the Knaster--Tarski fixpoint theorem.  The set
$\mathrm{Sec}(\mathcal{F})$ ordered by $\sqsubseteq$ forms a
complete lattice (it is a sub-lattice of the product
$\prod_{U \in \mathcal{U}} \mathcal{F}(U)$, which is complete because
each $\mathcal{F}(U)$ has a discrete complete order and the product
of complete lattices is complete).

Define the operator $\Phi : \mathrm{Sec}(\mathcal{F}) \to
\mathrm{Sec}(\mathcal{F})$ by
\[
  \Phi(\sigma) = \bigsqcup \{\sigma' \in \mathrm{Sec}(\mathcal{F})
  \mid h(\sigma') \geq h(\sigma)\}.
\]
We verify that $\Phi$ is monotone: if $\sigma_1 \sqsubseteq
\sigma_2$ (i.e., $h(\sigma_1) \leq h(\sigma_2)$), then
$\{\sigma' \mid h(\sigma') \geq h(\sigma_2)\} \subseteq
\{\sigma' \mid h(\sigma') \geq h(\sigma_1)\}$, so
$\Phi(\sigma_2) \sqsubseteq \Phi(\sigma_1)$.  (Note the reversal:
$\Phi$ is monotone on the \emph{opposite} lattice, which suffices for
Knaster--Tarski on the dual.)

By Knaster--Tarski, $\Phi$ has a greatest fixpoint
$\sigma^* = \bigsqcup\{\sigma \mid \sigma \sqsubseteq \Phi(\sigma)\}$.
This fixpoint satisfies $h(\sigma^*) = h(\Phi(\sigma^*))$, meaning
no further refinement improves the harmony: $\sigma^*$ is a
harmony equilibrium.
\end{proof}

\begin{proposition}[Uniqueness of harmony fixpoint]
\label{prop:harmony-unique}
If the harmony function $h$ is \emph{strictly monotone} (i.e., for
each theory $\mathcal{T}_k$ and each section $\sigma$, there is at
most one $\sigma'$ maximising $g_k$ among sections compatible
with $\sigma$), then the fixpoint is unique.
\end{proposition}

\begin{proof}
Suppose $\sigma_1^*$ and $\sigma_2^*$ are two fixpoints.
By strict monotonicity, $\sigma_1^* \sqsubseteq \sigma_2^*$ and
$\sigma_2^* \sqsubseteq \sigma_1^*$, giving
$h(\sigma_1^*) = h(\sigma_2^*)$.  Since each $g_k$ has a unique
maximiser, the sections must agree on each stratum, hence
$\sigma_1^* = \sigma_2^*$.
\end{proof}

%% ── §20.3  The Bottleneck Theorem ───────────────────────────────────
\section{The bottleneck theorem}
\label{sec:bottleneck}

\begin{theorem}[Bottleneck principle]
\label{thm:bottleneck}
Let $h(\sigma) = \bigotimes_{k=1}^{N} w_k \otG g_k(\sigma)$ be the
harmony function.  If $(\Grade, \otG)$ is an idempotent commutative
monoid with $\otG = \min$ (the bottleneck semiring), then:
\[
  h(\sigma) = \min_{k=1}^{N}\; w_k \otG g_k(\sigma),
\]
and the minimum is achieved by at least one stratum $k^*$.
\end{theorem}

\begin{proof}
When $\otG = \min$, the tensor product of a finite family is the
minimum of the family (by induction on $N$: for $N = 1$ the result is
trivial; for $N+1$, $\min(a_1, \ldots, a_N, a_{N+1}) =
\min(\min(a_1, \ldots, a_N), a_{N+1})$ by associativity of $\min$).

Since $\Grade$ is totally ordered and $\{w_k \otG g_k(\sigma)\}$ is
a finite set of elements of $\Grade$, the minimum is achieved by
some index $k^* \in \{1, \ldots, N\}$.  This is the
\emph{bottleneck stratum}: the stratum whose grade most constrains
the overall harmony.
\end{proof}

\begin{corollary}[Bottleneck-driven repair]
\label{cor:bottleneck-repair}
To improve $h(\sigma)$, it is necessary and sufficient to improve the
grade at the bottleneck stratum $k^*$.  Improving any other stratum
$k \neq k^*$ does not change $h(\sigma)$ (until the bottleneck
shifts).
\end{corollary}

\begin{proof}
Since $h(\sigma) = w_{k^*} \otG g_{k^*}(\sigma) =
\min_k(w_k \otG g_k(\sigma))$, increasing $g_k(\sigma)$ for
$k \neq k^*$ does not change the minimum (the minimum is still
$w_{k^*} \otG g_{k^*}(\sigma)$).  Increasing $g_{k^*}(\sigma)$ to
$g_{k^*}'$ raises the harmony to
$\min(w_{k^*} \otG g_{k^*}', \min_{k \neq k^*} w_k \otG g_k)$,
which is $> h(\sigma)$ provided $w_{k^*} \otG g_{k^*}' >
w_{k^*} \otG g_{k^*}$.
\end{proof}

%% ── §20.4  FMLM Convergence ────────────────────────────────────────
\section{$\FMLM$ convergence}
\label{sec:fmlm}

\begin{definition}[$\FMLM$ refinement operator]
\label{def:fmlm-operator}
The \emph{Fill-in by Masked Language Model} ($\FMLM$) procedure
operates on a partially specified section $\sigma_0$ (a ``template''
with masked slots).  The refinement operator
$R : \mathrm{Sec}(\mathcal{F}) \to \mathrm{Sec}(\mathcal{F})$ is:
\begin{enumerate}[label=(\arabic*),nosep]
\item \textbf{Identify bottleneck:} Compute $k^* =
  \arg\min_k\, w_k \otG g_k(\sigma_n)$.
\item \textbf{Generate candidates:} Produce a finite set
  $\mathrm{Cand}(\sigma_n, k^*)$ of candidate replacements for the
  masked slots in stratum $k^*$.
\item \textbf{Grade candidates:} For each candidate $\sigma' \in
  \mathrm{Cand}$, compute $h(\sigma')$.
\item \textbf{Select best:} Set $R(\sigma_n) = \arg\max_{\sigma'}
  h(\sigma')$.
\end{enumerate}
The iteration is $\sigma_{n+1} = R(\sigma_n)$.
\end{definition}

\begin{theorem}[$\FMLM$ convergence]
\label{thm:fmlm}
Let $(\Grade, \leq)$ be a complete lattice with the ascending chain
condition (ACC) (equivalently, $\Grade$ is finite or well-ordered).
Suppose the $\FMLM$ refinement operator $R$ satisfies:
\begin{enumerate}[label=(\roman*),nosep]
\item \emph{Non-degeneracy:} If $\sigma_n$ is not a fixpoint, then
  $\mathrm{Cand}(\sigma_n, k^*)$ contains at least one candidate
  $\sigma'$ with $h(\sigma') > h(\sigma_n)$.
\item \emph{Budget bound:} The total number of iterations is bounded
  by $B < \infty$.
\end{enumerate}
Then the sequence $(\sigma_n)_{n \geq 0}$ converges to a fixpoint
$\sigma^*$ with $h(\sigma^*) \geq h(\sigma_0)$, in at most
$\min(B, |\Grade|)$ steps.
\end{theorem}

\begin{proof}
The harmony sequence $(h(\sigma_n))_{n \geq 0}$ is non-decreasing:
$h(\sigma_{n+1}) = h(R(\sigma_n)) \geq h(\sigma_n)$ because $R$
selects $\arg\max h$ from a candidate set that includes $\sigma_n$
itself (the ``no change'' candidate).

By non-degeneracy, if $\sigma_n$ is not a fixpoint, then
$h(\sigma_{n+1}) > h(\sigma_n)$ (strict increase).

Since $(\Grade, \leq)$ satisfies ACC, every strictly ascending chain
is finite.  Thus the sequence $(h(\sigma_n))$ stabilises after
finitely many steps, say at step $n^* \leq |\Grade|$.  At $n^*$,
$h(\sigma_{n^*+1}) = h(\sigma_{n^*})$, so either $\sigma_{n^*}$ is
a fixpoint of $R$ (the non-degeneracy condition is not triggered)
or the iteration has exhausted the budget $B$.

With the budget bound, the iteration stops at
$n^* = \min(B, |\Grade|)$.

The inequality $h(\sigma^*) \geq h(\sigma_0)$ follows from the
non-decreasing property of the harmony sequence.
\end{proof}

\begin{remark}[Convergence rate]
\label{rem:convergence-rate}
If the grade semiring is $\Grade = [0,1]$ (discretised to $M$
levels), then $|\Grade| = M$ and convergence occurs in at most $M$
steps.  In practice, the bottleneck principle
(\cref{thm:bottleneck}) ensures that each step improves the
bottleneck stratum, and the number of strata $N$ provides a tighter
bound: at most $N \cdot M$ refinements suffice (each stratum can be
the bottleneck at most $M$ times before it is no longer minimal).
\end{remark}

%% ── §20.5  Descent as Optimisation ──────────────────────────────────
\section{Descent as optimisation}
\label{sec:descent-opt}

The final result of this chapter establishes that descent (finding a
global section of a fibration) is equivalent to optimisation (finding
the globally optimal section under the harmony metric).

\begin{theorem}[Descent-optimisation correspondence]
\label{thm:descent-opt}
Let $p : \mathcal{E} \to \mathcal{B}$ be a fibration of categories,
$\mathcal{U} = \{U_i \to X\}$ a cover in $\mathcal{B}$, and
$\mathcal{F}$ the associated presheaf of sections.  Let
$h : \mathrm{Sec}(\mathcal{F}) \to \Grade$ be the harmony function.
Then:
\begin{enumerate}[label=(\roman*)]
\item A \emph{descent datum} $(s_i, \phi_{ij})$ (local sections
  with gluing isomorphisms) corresponds to a \emph{candidate section}
  $\sigma \in \mathrm{Sec}(\mathcal{F})$.
\item Two descent data are \emph{coboundary equivalent}
  ($(s_i, \phi_{ij}) \sim (s_i', \phi_{ij}')$ iff there exist
  isomorphisms $\psi_i : s_i \xrightarrow{\sim} s_i'$ compatible with
  the $\phi$'s) if and only if the corresponding candidate sections
  are \emph{gauge equivalent}: $\sigma \sim_g \sigma'$ iff there
  exists $(t_i)$ with $\sigma' = \sigma + \delta^0 t$.
\item The \emph{cocycle condition} ($\phi_{ij} \circ \phi_{jk} =
  \phi_{ik}$ on triple overlaps) is equivalent to \emph{local
  optimality}: $\sigma$ is a critical point of $h$ restricted to
  each overlap region.
\item A descent datum is \emph{effective} (glues to a global section)
  if and only if the corresponding candidate section is \emph{globally
  optimal}: $h(\sigma) = \max_{\sigma'} h(\sigma')$.
\end{enumerate}
\end{theorem}

\begin{proof}
(i)\; A descent datum assigns to each $U_i$ a section
$s_i \in \mathcal{F}(U_i)$ and to each overlap $U_{ij} = U_i
\times_X U_j$ an isomorphism $\phi_{ij} : s_i|_{U_{ij}}
\xrightarrow{\sim} s_j|_{U_{ij}}$.  This is precisely a 0-cochain
$(s_i) \in \Cech^0(\mathcal{U}, \mathcal{F})$ together with a
1-cochain $(\phi_{ij}) \in \Cech^1(\mathcal{U},
\mathrm{Aut}(\mathcal{F}))$.  The candidate section
$\sigma = (s_i)_{i \in I}$ is the 0-cochain.

(ii)\; A coboundary equivalence is a 0-cochain $(\psi_i)$ with
$\phi_{ij}' = \psi_j \circ \phi_{ij} \circ \psi_i^{-1}$, which is
$(\phi') = (\phi) + \delta^0(\psi)$ in the \v{C}ech complex of
$\mathrm{Aut}(\mathcal{F})$.  On sections, $s_i' = \psi_i(s_i)$,
so $\sigma' = \sigma + \delta^0 t$ where $t_i$ is the displacement
$\psi_i(s_i) - s_i$.  This is gauge equivalence.

(iii)\; The cocycle condition $\phi_{ij} \circ \phi_{jk} = \phi_{ik}$
means the discrepancies $d_{ij} = s_i|_{U_{ij}} - s_j|_{U_{ij}}$
satisfy $d_{ij} + d_{jk} = d_{ik}$, i.e., $(d_{ij}) \in
\ker \delta^1$.  In terms of harmony: the grade $g(\sigma)$
restricted to the overlap $U_{ij}$ depends on the discrepancy
$d_{ij}$.  The cocycle condition says all overlap discrepancies are
coherent, which is local optimality: no local rearrangement of
sections on overlaps can improve the harmony (since the discrepancies
are already minimal---they satisfy the tightest possible constraint).

(iv)\; If the descent datum is effective, there exists a global
section $s \in \mathcal{F}(X)$ with $s|_{U_i} = s_i$.  This global
section has harmony
\[
  h(s) = \bigotimes_k w_k \otG g_k(s) = \gone
\]
(the maximal grade) because there is no discrepancy ($d_{ij} = 0$,
$H^1 = 0$).  Conversely, if $h(\sigma) = \max h$, then all overlap
discrepancies vanish (the bottleneck is $\gone$), the cocycle
condition is satisfied trivially, and the descent datum is effective
by the sheaf condition.
\end{proof}

\begin{corollary}[Equivalence as descent on product site]
\label{cor:equiv-descent}
Two artifacts $P$ and $Q$ are equivalent with respect to a property
$\varphi$ if and only if
$\check{H}^1(\mathcal{U}_{\mathrm{eq}},
\mathcal{F}_{\mathrm{eq}}) = 0$
on the product site $\mathcal{C}(P) \times \mathcal{C}(Q)$, where
$\mathcal{F}_{\mathrm{eq}}(U_P, U_Q) =
\{(s_P, s_Q) \mid s_P|_{U_P} \sim_\varphi s_Q|_{U_Q}\}$.
\end{corollary}

\begin{proof}
The equivalence $P \sim_\varphi Q$ means every local section
witnessing $\varphi$ on a region of $P$ corresponds to a local
section on the corresponding region of $Q$, and these correspondences
are compatible on overlaps.  This is precisely a descent datum on the
product site with the relational presheaf $\mathcal{F}_{\mathrm{eq}}$.
By \cref{thm:fundamental}, the descent datum is effective iff
$H^1 = 0$.
\end{proof}

\begin{proposition}[Support-local invalidation]
\label{prop:support-local}
If evidence changes are supported on a sub-region $V \subseteq X$,
then only judgments in
\[
  \mathrm{Star}_{\mathcal{U}}(V) = \{U_i \in \mathcal{U}
  \mid U_i \times_X V \neq \emptyset\}
\]
may need re-verification.
\end{proposition}

\begin{proof}
Let $\sigma$ be a section and $\sigma'$ the modified section with
changes supported on~$V$.  Then $\sigma|_{U_i} = \sigma'|_{U_i}$ for
all $U_i$ with $U_i \times_X V = \emptyset$.  The discrepancy
$d_{ij}' = \sigma_i'|_{U_{ij}} - \sigma_j'|_{U_{ij}}$ equals
$d_{ij}$ whenever both $U_i$ and $U_j$ are disjoint from $V$.  Thus
the cocycle condition is preserved on all overlaps outside
$\mathrm{Star}(V)$, and only overlaps involving at least one member
of $\mathrm{Star}(V)$ need re-checking.

By the Mayer--Vietoris principle, the cohomology of $X$ decomposes:
the restriction map $H^1(X) \to H^1(X \setminus V)$ is an isomorphism
on the complement of $\mathrm{Star}(V)$, so only the contribution
from $\mathrm{Star}(V)$ can change.
\end{proof}

%%%%%%%%%%%%%%%%%%%%%%%%%%%%%%%%%%%%%%%%%%%%%%%%%%%%%%%%%%%%%%%%%%%%%%%%
%%  Chapter 21 — Metatheoretic Properties
%%%%%%%%%%%%%%%%%%%%%%%%%%%%%%%%%%%%%%%%%%%%%%%%%%%%%%%%%%%%%%%%%%%%%%%%
\chapter{Metatheoretic Properties}
\label{ch:metatheory}

This chapter establishes consistency, conservativity, decidability
results, completeness of the subsumption, and cut elimination for
the graded type system.

%% ── §21.1  Consistency ─────────────────────────────────────────────
\section{Consistency}
\label{sec:consistency}

\begin{theorem}[Consistency of the 4-tuple framework]
\label{thm:consistency}
Assuming ZFC with a Grothendieck universe $\mathbb{U}$, the 4-tuple
framework $(\STop, \PStr, \CSpec, \GEval)$ is consistent: there is no
derivation of $\Gamma \vdash_g E : \bot$ for any context $\Gamma$,
grade $g$, and term $E$.
\end{theorem}

\begin{proof}
We construct a model of the 4-tuple framework in the topos
$\Set_{\mathbb{U}}$ of $\mathbb{U}$-small sets.

\emph{Step 1: Model of $\STop$.}
Take $\mathcal{C} = \mathbf{1}$ (the terminal category) with the
maximal topology.  Then $\Sh(\mathbf{1}) \cong \Set_{\mathbb{U}}$,
which is an elementary topos with natural number object.  Set
$\mathcal{F} = \mathrm{const}_{\{*\}}$ (the constant presheaf on a
singleton).  The subobject classifier is $\Omega = \{0, 1\}$, the
sheafification is the identity, and the global-sections functor
$\gamma^*$ is the identity.  This is a valid Semantic Topos.

\emph{Step 2: Model of $\PStr$.}
Take $\Gamma = ()$ (empty context), $A = \mathbb{1}$ (the unit
type), and $E = * : \mathbb{1}$ (the unique inhabitant).  The type
theory is standard MLTT interpreted in $\Set_{\mathbb{U}}$ via the
standard set-theoretic semantics \cite{mlm1984}.  Under this
interpretation:
\begin{itemize}[nosep]
\item $\Pi$-types are interpreted as dependent products (indexed
  families of functions);
\item $\Sigma$-types as dependent sums (indexed disjoint unions);
\item Identity types $\mathrm{Id}_A(a,b)$ as the singleton
  $\{*\}$ if $a = b$ and $\emptyset$ otherwise;
\item The universe $\Type_0$ is interpreted as $\mathbb{U}$ itself.
\end{itemize}
The empty type $\bot$ is interpreted as $\emptyset$.  Since there is
no set-theoretic function $f : \Gamma \to \emptyset$ when $\Gamma$
is non-empty (or when $\Gamma$ is the terminal object), there is no
closed term of type $\bot$ in the model.

\emph{Step 3: Model of $\CSpec$.}
The trivial cover $\{* \to *\}$ on $\mathbf{1}$ gives
$\Cech^0 = \{*\}$ and $\Cech^n = 0$ for $n \geq 1$.  All cohomology
vanishes: $H^n = 0$.  This is a valid Cohomological Spectrum.

\emph{Step 4: Model of $\GEval$.}
Take $G = \mathbf{2} = \{\gzero, \gone\}$ with $\otG = \land$,
$\opG = \lor$.  The evaluation $\mu(\mathcal{D}) = \gone$ for every
valid derivation.  This is a valid quantale.

\emph{Step 5: Compatibility.}
The compatibility conditions (K1)--(K4) are satisfied trivially in
this model: the trivial cover refines the maximal topology (K1),
the presheaf is shared (K2), the unit type classifies sections
of the constant presheaf (K3), and the evaluation is well-defined
(K4).

\emph{Step 6: Consistency.}
If $\bot$ were derivable in the 4-tuple framework, there would exist
$\Gamma, g, E$ with $\Gamma \vdash_g E : \bot$.  Under the
set-theoretic interpretation, this would give an element of $\emptyset$,
which is impossible.  Hence the framework is consistent.
\end{proof}

%% ── §21.2  Conservativity ──────────────────────────────────────────
\section{Conservativity over MLTT}
\label{sec:conservativity}

\begin{theorem}[Conservative extension]
\label{thm:conservative}
Adding the Grade semiring $\Grade$ to Martin-L\"of Type Theory is
conservative over MLTT: if $\Gamma \vdash_{\gone} E : A$ is
derivable in the graded system and $\Gamma, A, E$ contain no grade
annotations, then $\Gamma \vdash E : A$ is derivable in standard
MLTT.
\end{theorem}

\begin{proof}
We define a \emph{grade-erasure} map $|-| : \mathrm{Judg}_g \to
\mathrm{Judg}$ from graded judgments to ordinary MLTT judgments,
and show it preserves derivability for grade-$\gone$ judgments.

\emph{Grade erasure.}  Define $|-|$ on judgments, types, and terms by:
\begin{align*}
  |\Gamma \vdash_g E : A| &= \Gamma \vdash |E| : |A|, \\
  |\GId_g(A, a, b)| &= \mathrm{Id}_{|A|}(|a|, |b|)
    \quad\text{when } g = \gone, \\
  |\Pi_g(x : A).\, B| &= \Pi(x : |A|).\, |B|
    \quad\text{when } g = \gone, \\
  |\Sigma_g(x : A).\, B| &= \Sigma(x : |A|).\, |B|
    \quad\text{when } g = \gone.
\end{align*}
For $g < \gone$, the graded types have no MLTT counterpart, but we
are only considering $g = \gone$ judgments.

\emph{Preservation of derivations.}  We show by induction on the
derivation of $\Gamma \vdash_{\gone} E : A$ that
$\Gamma \vdash |E| : |A|$ is derivable in MLTT.

\textsc{Case Variable:}
$\dfrac{(x : A) \in \Gamma}{\Gamma \vdash_{\gone} x : A}$.
Erasing: $(x : |A|) \in \Gamma$, so $\Gamma \vdash x : |A|$ by the
MLTT variable rule.

\textsc{Case $\Pi$-introduction:}
$\dfrac{\Gamma, x : A \vdash_{\gone} E : B}
  {\Gamma \vdash_{\gone} \lambda x.\, E : \Pi_{\gone}(x:A).\, B}$.
By induction, $\Gamma, x : |A| \vdash |E| : |B|$.  By the MLTT
$\Pi$-introduction rule,
$\Gamma \vdash \lambda x.\, |E| : \Pi(x:|A|).\, |B|$.

\textsc{Case $\Pi$-elimination:}
$\dfrac{\Gamma \vdash_{\gone} f : \Pi_{\gone}(x:A).\, B \quad
  \Gamma \vdash_{\gone} a : A}
  {\Gamma \vdash_{\gone} f\,a : B[a/x]}$.
By induction, $\Gamma \vdash |f| : \Pi(x:|A|).\,|B|$ and
$\Gamma \vdash |a| : |A|$.  By MLTT $\Pi$-elimination,
$\Gamma \vdash |f|\,|a| : |B|[|a|/x]$.

\textsc{Case Grade-weakening:}
$\dfrac{\Gamma \vdash_g E : A \quad g \leq g'}
  {\Gamma \vdash_{g'} E : A}$.
When $g' = \gone$, we need $g = \gone$ (since $\gone$ is maximal
and $g \leq \gone$ always holds; but if $g < \gone$, the premise
has grade $< \gone$ and the erasure may not apply).  When
$g = \gone$, the premise directly gives the result by induction.

The remaining cases ($\Sigma$-types, identity types, universe rules,
$W$-types) are analogous: at grade $\gone$, every graded type former
reduces to its MLTT counterpart, and every graded rule reduces to
the corresponding MLTT rule.

\emph{Conservativity.}  If $\Gamma \vdash_{\gone} E : A$ is
derivable in the graded system with $\Gamma, A, E$ containing no
grade annotations, then the grade-erasure $|\Gamma \vdash_{\gone}
E : A| = \Gamma \vdash E : A$ is derivable in MLTT by the induction
above.  Hence the graded system proves no new $\gone$-level judgments
about pure MLTT types and terms.
\end{proof}

%% ── §21.3  Decidability ────────────────────────────────────────────
\section{Decidability results}
\label{sec:decidability}

\begin{theorem}[Decidability of grade computation]
\label{thm:dec-grade}
For a finite number of strata $N$ and a decidable grade semiring
$\Grade$, the harmony computation
$h(\sigma) = \bigotimes_{k=1}^{N} w_k \otG g_k(\sigma)$
is decidable.
\end{theorem}

\begin{proof}
The computation requires:
\begin{enumerate}[label=(\arabic*),nosep]
\item Evaluating each $g_k(\sigma)$ for $k = 1, \ldots, N$.  Since
  $\Grade$ is decidable and each $g_k$ is a computable function (by
  hypothesis), each evaluation terminates.
\item Computing $N - 1$ applications of $\otG$.  Since $\otG$ is a
  computable operation on a decidable semiring, each application
  terminates.
\item Computing $N$ applications of the weight multiplication
  $w_k \otG (-)$.  Again computable.
\end{enumerate}
The total computation is a finite sequence of decidable operations,
hence decidable.  The complexity is $O(N \cdot C_g)$ where $C_g$ is
the cost of a single grade evaluation.
\end{proof}

\begin{theorem}[Decidability of harmony]
\label{thm:dec-harmony}
If each theory slot $\mathcal{T}_k$ has a decidable satisfaction
relation (i.e., given $\sigma$ and $\mathcal{T}_k$, it is decidable
whether $g_k(\sigma) \geq t$ for any threshold $t$), then the harmony
comparison $h(\sigma_1) \geq h(\sigma_2)$ is decidable.
\end{theorem}

\begin{proof}
By \cref{thm:dec-grade}, $h(\sigma_1)$ and $h(\sigma_2)$ are both
computable.  The comparison $h(\sigma_1) \geq h(\sigma_2)$ reduces to
$\leq$ on $\Grade$, which is decidable by hypothesis.
\end{proof}

\begin{theorem}[Decidability of graded type-checking]
\label{thm:dec-typecheck}
If the underlying type theory has decidable type-checking (given
$\Gamma, E, A$, it is decidable whether $\Gamma \vdash E : A$),
then graded type-checking ($\Gamma \vdash_g E : A$) is decidable.
\end{theorem}

\begin{proof}
Graded type-checking extends ordinary type-checking with grade
tracking.  Each typing rule of the graded system has the form:
\[
  \dfrac{\Gamma \vdash_{g_1} E_1 : A_1 \quad \cdots \quad
    \Gamma \vdash_{g_n} E_n : A_n \quad g = f(g_1, \ldots, g_n)}
    {\Gamma \vdash_g E : A}
\]
where $f$ is a computable function of the input grades (typically
$\otG$ or $\opG$).  Given a purported derivation, checking each rule
application requires:
\begin{enumerate}[label=(\arabic*),nosep]
\item Checking the underlying MLTT judgment (decidable by hypothesis).
\item Computing the grade from the sub-grades (decidable by
  \cref{thm:dec-grade}).
\item Comparing the computed grade with the asserted grade (decidable
  on $\Grade$).
\end{enumerate}
By induction on the derivation tree, the full check is decidable.
\end{proof}

\begin{theorem}[Undecidability of full descent]
\label{thm:undec-descent}
The descent problem---``given a site $(\mathcal{C}, \mathcal{J})$, a
presheaf $\mathcal{F}$, and a cover $\mathcal{U}$, does
$\check{H}^1(\mathcal{U}, \mathcal{F}) = 0$?''---is undecidable in
general.
\end{theorem}

\begin{proof}
We reduce the halting problem to descent.  Given a Turing machine $M$
and input $w$, construct a site $\mathcal{C}_M$ as follows:
\begin{itemize}[nosep]
\item Objects: the configurations $(q, \tau, h)$ of $M$ on input~$w$
  (state, tape, head position), plus a terminal object~$X$.
\item Morphisms: one-step transitions $c \to c'$ and the unique
  morphism $c \to X$ for each configuration.
\item Cover: $\mathcal{U} = \{c \to X \mid c \text{ reachable}\}$.
\end{itemize}
Define the presheaf $\mathcal{F}(c) = \{*\}$ if $c$ is not a
halting configuration, and $\mathcal{F}(c) = \emptyset$ if $c$ is a
halting configuration.

If $M$ halts on $w$, some configuration $c_h$ has
$\mathcal{F}(c_h) = \emptyset$, and the local sections on
$c_h$'s neighbors cannot be extended through $c_h$.  This creates a
non-trivial 1-cocycle: $\check{H}^1 \neq 0$.

If $M$ does not halt, every configuration has
$\mathcal{F}(c) = \{*\}$, all restrictions are identities, and the
constant family $(*)$ is a compatible descent datum:
$\check{H}^1 = 0$.

Thus $\check{H}^1 = 0$ iff $M$ does not halt on $w$, which is
undecidable.
\end{proof}

\begin{theorem}[Decidability for bounded depth]
\label{thm:dec-bounded}
If the site $\mathcal{C}$ has finite depth $d$ (every chain of
morphisms has length $\leq d$) and $\mathcal{F}$ takes values in
finite sets, then descent is decidable.
\end{theorem}

\begin{proof}
Under these finiteness conditions:
\begin{enumerate}[label=(\arabic*),nosep]
\item The \v{C}ech complex $\Cech^\bullet(\mathcal{U}, \mathcal{F})$
  is a finite complex of finite-dimensional vector spaces (over the
  appropriate coefficient ring).
\item Each coboundary map $\delta^n$ is a finite matrix.
\item $\ker \delta^n$ and $\mathrm{im}\, \delta^{n-1}$ are computable
  by Gaussian elimination (or its analogue for abelian groups).
\item $\check{H}^1 = \ker \delta^1 / \mathrm{im}\, \delta^0$ is
  therefore computable.
\item Checking $\check{H}^1 = 0$ reduces to checking whether
  $\ker \delta^1 = \mathrm{im}\, \delta^0$, which is decidable.
\end{enumerate}
The complexity is polynomial in the size of the \v{C}ech complex:
$O(|\mathcal{U}|^{d+1} \cdot |\mathcal{F}|^3)$ for the Gaussian
elimination step.
\end{proof}

%% ── §21.4  Completeness of the Subsumption ──────────────────────────
\section{Completeness of the subsumption}
\label{sec:completeness}

\begin{theorem}[Proper extension --- detailed version]
\label{thm:proper-detail}
The JuGeo 4-tuple framework is a proper extension of every subsumed
framework: it is not equivalent to any of Montague Grammar, DRT/SDRT,
MLTT, OT, Linear Logic, the Generative Lexicon, HoTT, or Prior's
Tense Logic.
\end{theorem}

\begin{proof}
By \cref{thm:proper}, the separation example $\JJ^*$ lies outside
every embedding.  We sharpen this by exhibiting a judgment that
\emph{requires} all four components simultaneously.

\emph{The separation judgment $\JJ^*$.}
Consider the judgment ``The poem $P$ achieves a unity of meter,
imagery, and meaning across its three stanzas.''  Formalise this as:

\begin{itemize}[nosep]
\item $\STop^*$: Site $\mathcal{C}^*$ with objects
  $\{s_1, s_2, s_3, P\}$ (three stanzas and the full poem), covers
  $\{s_1, s_2, s_3\} \to P$ (the stanzas cover the poem).
  Three strata: $\phon$ (meter), $\sem$ (imagery), $\prag$
  (meaning).  The topos $\Sh(\mathcal{C}^*)$ is non-trivial (not
  equivalent to $\Set$).

\item $\PStr^*$: The graded judgment
  $\Gamma \vdash_{0.7} E_P : \mathrm{Unity}(P)$
  where $\mathrm{Unity}(P)$ is a dependent type encoding the three
  strata of unity.  The grade $0.7$ (not binary) records that the
  meter is excellent ($g_{\phon} = 0.95$), the imagery is strong
  ($g_{\sem} = 0.85$), but the meaning has a forced metaphor
  ($g_{\prag} = 0.7$).

\item $\CSpec^*$: The \v{C}ech complex on the stanza cover yields
  $H^1 \neq 0$: the metaphor in the transition from $s_2$ to $s_3$
  creates a cohomological obstruction (the local meanings do not glue
  smoothly).

\item $\GEval^*$: $G = [0,1]$ with $\otG = \min$.
  $h(\JJ^*) = \min(0.95, 0.85, 0.7) = 0.7$ (the meaning stratum
  is the bottleneck).
\end{itemize}

This judgment simultaneously requires:
\begin{enumerate}[label=(\roman*),nosep]
\item A non-trivial site (rules out MLTT, HoTT, Linear Logic,
  Generative Lexicon);
\item Non-binary grades (rules out Montague, DRT, MLTT, HoTT, Tense
  Logic);
\item Non-trivial $H^1$ (rules out all frameworks that have only
  trivial cohomology);
\item Multi-stratal harmony (rules out OT, which has a single stratum
  with ranked constraints but no cross-stratal harmony).
\end{enumerate}

No single subsumed framework can express all four properties.  Hence
JuGeo is a proper extension.
\end{proof}

%% ── §21.5  Cut Elimination ──────────────────────────────────────────
\section{Cut elimination for the graded system}
\label{sec:cut-elim}

\begin{definition}[Graded sequent calculus]
\label{def:graded-sequent}
The graded sequent calculus has judgments of the form
$\Gamma \vdash_g E : A$ where $\Gamma$ is a context, $g \in \Grade$
is a grade, $E$ is a term, and $A$ is a type.  The \emph{cut rule}
is:
\[
  \dfrac{\Gamma \vdash_g E : A \quad \Gamma, x : A \vdash_{g'} E' : B}
    {\Gamma \vdash_{g \otG g'} E'[E/x] : B}
  \;\;\textsc{(Cut)}
\]
The grade of the conclusion is $g \otG g'$: the tensor product of
the grades of the two premises.
\end{definition}

\begin{definition}[Cut-free derivation]
A derivation is \emph{cut-free} if it does not use the \textsc{(Cut)}
rule.
\end{definition}

\begin{theorem}[Cut elimination]
\label{thm:cut-elim}
If $\Gamma \vdash_g E : A$ has a derivation with cuts, then there
exists a cut-free derivation $\Gamma \vdash_{g'} E' : A$ with
$g' \geq g$ (the grade may increase, but does not decrease).
Moreover, $E' =_\beta E$ (the terms are $\beta$-equal).
\end{theorem}

\begin{proof}
We adapt Gentzen's cut elimination procedure, tracking grades through
each reduction step.

\emph{Measure.}  Define the \emph{cut-rank} of a derivation $\mathcal{D}$
as the pair $(\rho, \nu)$ where $\rho$ is the maximum complexity
(number of connectives) of a cut formula in $\mathcal{D}$, and $\nu$
is the number of cuts of that maximum complexity.  We order cut-ranks
lexicographically.

\emph{Key reduction.}  Consider a topmost cut (one whose premises
are cut-free or contain only lower-rank cuts):
\[
  \dfrac{\dfrac{\vdots}{\Gamma \vdash_{g_1} E_1 : A}
    \quad
    \dfrac{\vdots}{\Gamma, x:A \vdash_{g_2} E_2 : B}}
    {\Gamma \vdash_{g_1 \otG g_2} E_2[E_1/x] : B}
  \;\;\textsc{(Cut)}
\]

\textsc{Case 1: Principal cut on $\Pi$-types.}
The left premise ends with $\Pi$-introduction and the right with
$\Pi$-elimination:
\begin{gather*}
  \dfrac{\Gamma, y:C \vdash_{g_1} E_1 : D}
    {\Gamma \vdash_{g_1} \lambda y.\,E_1 : \Pi(y:C).\,D}
  \quad
  \dfrac{\Gamma \vdash_{g_2} f : \Pi(y:C).\,D \quad
    \Gamma \vdash_{g_3} a : C}
    {\Gamma \vdash_{g_2 \otG g_3} f\,a : D[a/y]}
\end{gather*}
Cutting on $\Pi(y:C).\,D$: the cut reduces to substituting $\lambda
y.\,E_1$ for $f$, giving $(\lambda y.\,E_1)\,a \to_\beta E_1[a/y]$.
The grade of the result is $g_1 \otG g_3$: we lose the $g_2$ factor
(the grade of the eliminated function) but gain back $g_1$ (the
grade of the introduction).  Since $g_2 \leq \gone$ (grades are
bounded by $\gone$) and $g_1 \otG g_3 \geq g_1 \otG g_2 \otG g_3$
(because $g_2 \leq \gone$ and $\otG$ is monotone), the resulting
grade satisfies $g' = g_1 \otG g_3 \geq g_1 \otG g_2 \otG g_3 = g$.

The cut on $\Pi(y:C).\,D$ is replaced by cuts on $C$ and $D$, both
of lower complexity.  Cut-rank decreases.

\textsc{Case 2: Principal cut on $\Sigma$-types.}
Similar: $\Sigma$-introduction paired with projection eliminates the
cut, replacing it with lower-complexity cuts.  The grade bound
$g' \geq g$ holds by monotonicity of $\otG$.

\textsc{Case 3: Principal cut on $\mathrm{Id}$-types.}
The left premise introduces a reflexivity proof
$\mathrm{refl}_a : \mathrm{Id}_A(a,a)$ and the right eliminates it
via $J$ (path induction).  The reduction applies the $J$-computation
rule: $J(\mathrm{refl}_a, d) \to_\beta d$.  The grade computation
$g' = g_1 \otG g_d \geq g_1 \otG g_2 \otG g_d$ follows from
$g_2 \leq \gone$.

\textsc{Case 4: Non-principal cut (commutative case).}
If the cut formula is not the principal formula of one of the
premises, commute the cut upward past the last rule of that premise.
This does not change the cut-rank but moves the cut closer to an
axiom, where it can be eliminated trivially.

\emph{Grade tracking.}  At each reduction step:
\begin{itemize}[nosep]
\item Principal cuts: the reduced derivation has grade $g' \geq g$
  because the eliminated intermediate formula had grade $\leq \gone$,
  and removing it can only increase the tensor product (since
  $a \otG b \leq a$ when $b \leq \gone$ for the bottleneck semiring,
  so removing the $b$ factor gives $a \geq a \otG b$).
\item Commutative cuts: grade is unchanged (the cut is merely
  repositioned).
\end{itemize}

\emph{Termination.}  Each principal-cut reduction strictly decreases
the cut-rank $(\rho, \nu)$: either $\rho$ decreases (the maximum
cut complexity drops) or $\rho$ stays the same and $\nu$ decreases
(the number of maximal-complexity cuts drops by one).  Since
cut-ranks are well-ordered (pairs of natural numbers under
lexicographic order), the process terminates.

\emph{Conclusion.}  After finitely many reduction steps, all cuts
are eliminated.  The resulting derivation
$\Gamma \vdash_{g'} E' : A$ is cut-free, with $g' \geq g$ and
$E' =_\beta E$ (each reduction step is a $\beta$-reduction or a
commutation that preserves $\beta$-equality).
\end{proof}

\begin{corollary}[Subformula property]
\label{cor:subformula}
Every cut-free derivation in the graded sequent calculus has the
subformula property: every type appearing in the derivation is a
subformula of a type in the conclusion or the context.
\end{corollary}

\begin{proof}
The cut rule is the only rule that introduces a type ($A$, the cut
formula) not present in the conclusion.  In a cut-free derivation,
every type is introduced by an introduction or elimination rule,
which only involves subformulas of the conclusion type or context
types.
\end{proof}

\begin{corollary}[Consistency from cut elimination]
\label{cor:consistency-cut}
The graded sequent calculus is consistent: $\vdash_g \bot$ is not
derivable for any $g$.
\end{corollary}

\begin{proof}
By \cref{thm:cut-elim}, if $\vdash_g E : \bot$ were derivable, there
would be a cut-free derivation.  By the subformula property, every
type in the derivation is a subtype of $\bot$ or of a context type.
Since the context is empty and $\bot$ has no introduction rule, no
cut-free derivation of $\bot$ exists.
\end{proof}

%% ── §21.6  Structural Rules ────────────────────────────────────────
\section{Structural rules}
\label{sec:structural}

\begin{theorem}[Admissibility of structural rules]
\label{thm:structural}
The graded sequent calculus admits weakening, exchange, and
contraction, with the following grade effects:
\begin{enumerate}[label=(\roman*)]
\item \textbf{Weakening:} If $\Gamma \vdash_g E : B$, then
  $\Gamma, x : A \vdash_g E : B$ (grade unchanged).
\item \textbf{Exchange:} If $\Gamma, x : A, y : B, \Delta \vdash_g
  E : C$, then $\Gamma, y : B, x : A, \Delta \vdash_g E : C$
  (grade unchanged).
\item \textbf{Contraction:} If $\Gamma, x : A, y : A \vdash_g
  E : B$, then $\Gamma, z : A \vdash_{g'} E[z/x, z/y] : B$ with
  $g' \geq g$ (grade may increase).
\end{enumerate}
\end{theorem}

\begin{proof}
(i) Weakening:  The additional variable $x : A$ is unused in
the derivation of $E : B$.  Each typing rule in the derivation
remains valid with the extended context, since all rules are
monotone in the context.  The grade is unchanged because $x$
contributes no evidence.

(ii) Exchange: The typing rules depend only on the \emph{set} of
context bindings, not their order (contexts are multisets of
typed variables).

(iii) Contraction: Replacing two copies $x, y$ of a variable of
type $A$ with a single copy $z$ identifies two evidence paths.
If $x$ contributed grade $g_x$ and $y$ contributed grade $g_y$,
the contraction replaces $g_x \otG g_y$ with $g_z = g_x \opG g_y
\geq g_x \otG g_y$ (since $a \opG b \geq a \otG b$ in any
semiring where $\opG$ dominates $\otG$, which holds for the
idempotent semiring by the absorption law).
\end{proof}

%% ── §21.7  Source-Independent Soundness ─────────────────────────────
\section{Source-independent soundness}
\label{sec:source-sound}

\begin{theorem}[Source-independent soundness]
\label{thm:source-sound}
If descent succeeds on a family of local sections
$\{s_i\}_{i \in I}$, the resulting global section $s$ is valid
regardless of the source of each $s_i$ (whether produced by an SMT
solver, a runtime test, an LLM, or a human).  The grade of $s$ is
\[
  g(s) = \bigotimes_{i \in I} g(s_i),
\]
the tensor product of the grades of the individual sources.
\end{theorem}

\begin{proof}
The descent theorem (\cref{thm:fundamental}) is a purely mathematical
statement about the \v{C}ech complex: it depends only on the
\emph{values} of the local sections $s_i \in \mathcal{F}(U_i)$ and
their compatibility on overlaps, not on the \emph{process} that
produced them.

Concretely: the gluing condition
$s_i|_{U_{ij}} = s_j|_{U_{ij}}$ is a set-theoretic equality,
verifiable without reference to provenance.  If the equality holds,
the sheaf condition produces a unique global section $s$,
regardless of source.

The grade computation $g(s) = \bigotimes_i g(s_i)$ follows from the
multiplicativity of the grade evaluation: the grade of a composite
derivation is the tensor product of the grades of its parts
(by the definition of $\mu$ on derivation trees).  Each $g(s_i)$
reflects the trust level of its source (e.g., $g = \gone$ for
formal proofs, $g = 0.3$ for LLM suggestions), and the tensor
product computes the overall trust via the bottleneck principle.
\end{proof}

\begin{remark}
Source-independent soundness is the key property that enables
hybrid verification: combining formal proofs (high grade) with
LLM-generated suggestions (low grade) and runtime evidence (medium
grade) into a single judgment whose overall trust is the
conservative combination of all sources.
\end{remark}
